\documentclass[11pt,a4paper]{article}

\usepackage[polish]{babel}
\usepackage{fontspec}
\usepackage[a4paper,margin=2.5cm]{geometry}
\usepackage{amsmath,amssymb}
\usepackage{booktabs,longtable,tabularx,multirow,makecell}
\usepackage{graphicx}
\usepackage{float}
\usepackage{caption}
\usepackage{xcolor}
\usepackage[numbers,sort&compress]{natbib}
\usepackage[hidelinks]{hyperref}
\usepackage{enumitem}

\setlist{itemsep=2pt,topsep=4pt}
\captionsetup{font=small,labelfont=bf}
\captionsetup[table]{position=top}
\captionsetup[figure]{position=bottom}
\graphicspath{{figures/}}
\renewcommand{\arraystretch}{1.12}
\setlength{\emergencystretch}{2em}
\newcommand{\pp}{\,p.p.}

\title{Analiza miar różnorodności próbek w~baggingu\\[4pt]
\large Raport końcowy, projekt 1, Uczenie Zespołowe}
\author{Jerzy Czarnecki (327\,239) \and Małgorzata Dżaman (327\,241)}
\date{Politechnika Warszawska, Wydział Matematyki i~Nauk Informacyjnych\\
Warszawa, 24.11.2026}

\begin{document}

\maketitle

\begin{abstract}
\noindent
W~baggingu i~lesie losowym każde drzewo decyzyjne uczy się na osobnej losowej próbce danych treningowych, a~jakość zespołu zależy od dokładności drzew i~od tego, jak bardzo różnią się ich predykcje. Obie wielkości mierzy się zwykle dopiero po uczeniu. W~pracy sprawdziliśmy, czy proste miary różnorodności i~reprezentatywności samych próbek, obliczone przed uczeniem, są związane z~jakością zespołu. Porównaliśmy 15 sposobów losowania próbek: losowanie ze zwracaniem i~bez zwracania o~różnym rozmiarze, losowanie próbek o~wspólnej części przy stałym rozmiarze oraz losowanie warstwowe. Badanie objęło 12 zbiorów klasyfikacyjnych i~12 regresyjnych, pięciokrotną walidację krzyżową powtórzoną 10 razy i~łącznie 29\,400 zespołów po 100 drzew. Jednostką w~testach statystycznych był zbiór danych. Żaden sposób losowania nie dał lepszych zespołów niż losowanie stosowane domyślnie, czyli losowanie ze zwracaniem próbki o~rozmiarze części treningowej. Istotnie gorsze były próbki w~dużej części wspólne dla wszystkich drzew: przy tym samym rozmiarze próbek losowanie niezależne dało lepszy zespół w~21 z~24 zbiorów. Jakość zespołu zależała od liczby różnych obserwacji w~próbce, a~nie od tego, czy losowano ze zwracaniem. Miary próbek były istotnie, ale bardzo słabo związane z~różnicami między drzewami (współczynnik korelacji rang około 0.02--0.03). Po uwzględnieniu dokładności drzew większa różnorodność próbek wiązała się z~lepszym zespołem. Żadna pojedyncza miara nie wskazała jednak najlepszego sposobu losowania ani lepszego z~kilku losowań tym samym sposobem. Wyjaśniamy to tym, że przy ustalonym sposobie losowania miary próbek prawie się nie zmieniają, a~ich rozrzut maleje wraz z~liczbą obserwacji.

\end{abstract}

\tableofcontents
\newpage

% ====================================================================
\section{Cel i~problem badawczy}
\label{sec:cel}
% ====================================================================

Bagging~\cite{breiman1996} jest metodą budowy zespołu modeli. Z~części treningowej losuje się $M$ próbek, na każdej z~nich uczy się osobny model (w~tym projekcie drzewo decyzyjne), a~predykcje modeli uśrednia się. Las losowy~\cite{breiman2001} działa tak samo, ale dodatkowo w~każdym węźle drzewa losuje podzbiór cech, spośród których wybierany jest podział. W~obu metodach każdą próbkę losuje się domyślnie ze zwracaniem i~ma ona tyle obserwacji, ile liczy część treningowa. Takie losowanie nazywamy bootstrapowym.

Celem projektu jest sprawdzenie, czy miary różnorodności i~reprezentatywności próbek uczących, obliczone \emph{przed} uczeniem drzew, są związane z~jakością zespołu i~czy pozwalają wskazać sposób losowania próbek, przy którym zespół jest lepszy.

Breiman~\cite{breiman2001} pokazał, że błąd lasu drzew jest ograniczony z~góry przez
\[
PE^{*} \le \frac{\bar{\rho}\,(1-s^{2})}{s^{2}},
\]
gdzie $s$ oznacza siłę (dokładność) drzew, a~$\bar{\rho}$ ich średnią korelację. Obie wielkości zależą od sposobu losowania próbek: mniejsze próbki dają drzewa bardziej różne, ale słabsze. Znane miary różnorodności zespołu, takie jak statystyka Q, niezgodność czy korelacja błędów, wymagają predykcji nauczonych modeli~\cite{kuncheva2003,brown2005}. Próbki są natomiast znane przed uczeniem. Jeżeli miary obliczone na samych próbkach niosą informację o~jakości zespołu, można by na ich podstawie wybierać sposób losowania bez uczenia wielu zespołów.

\paragraph{Główny problem badawczy.} Czy miary różnorodności i~reprezentatywności próbek uczących, obliczone przed uczeniem modeli, są związane z~jakością zespołu, także po uwzględnieniu siły modeli składowych?

\paragraph{Pytania badawcze i~hipotezy.} Wszystkie hipotezy testujemy dwustronnie.
\begin{description}[leftmargin=1.2cm,style=nextline]
  \item[PB1] Czy różnice między próbkami przekładają się na różnice między drzewami?\\
  \textbf{H1:} drzewa uczone na bardziej różniących się próbkach dają bardziej różniące się predykcje, a~drzewa uczone na bardziej reprezentatywnych próbkach są dokładniejsze.
  \item[PB2] Czy miary próbek są związane z~jakością zespołu, także po uwzględnieniu średniej jakości drzew?\\
  \textbf{H2:} tak; w~szczególności przy stałym rozmiarze próbek większe pokrywanie się próbek różnych drzew pogarsza jakość zespołu.
  \item[PB3] Czy któryś sposób losowania daje lepsze zespoły niż losowanie bootstrapowe i~czy wskazują na to miary obliczone przed uczeniem?\\
  \textbf{H3:} co najmniej jeden wariant losowania daje zespół o~innej jakości niż losowanie bootstrapowe.
\end{description}

Nie zakładamy z~góry, że większa różnorodność próbek prowadzi do lepszego zespołu ani że jedna miara jest najlepsza dla wszystkich zbiorów danych.

% ====================================================================
\section{Przegląd literatury}
% ====================================================================

\paragraph{Bagging i~warianty próbkowania.} Bagging~\cite{breiman1996} zmniejsza wariancję niestabilnych modeli, takich jak głębokie drzewa decyzyjne. Breiman zaproponował też uczenie na małych fragmentach danych losowanych bez zwracania (\emph{pasting})~\cite{breiman1999}. Bühlmann i~Yu~\cite{buhlmann2002} wyjaśnili, w~jaki sposób bagging zmniejsza wariancję, i~zaproponowali \emph{subagging}, czyli losowanie mniejszych podpróbek bez zwracania. Pokazali też, że podpróbka o~rozmiarze $n/2$ daje w~przybliżeniu ten sam efekt co próbka bootstrapowa. Martínez-Muñoz i~Suárez~\cite{martinez2010} pokazali, że próbki mniejsze od części treningowej często dają lepsze zespoły, a~najlepszy rozmiar zależy od zbioru. Las losowy~\cite{breiman2001} dodaje do baggingu losowanie cech w~każdym węźle, co dodatkowo zmniejsza korelację drzew. Według Probsta i~in.~\cite{probst2019} rozmiar próbki i~sposób jej losowania (ze zwracaniem lub bez) są hiperparametrami lasu losowego, które wpływają na jego jakość, choć zwykle w~niewielkim stopniu.

\paragraph{Siła i~niezależność modeli.} Krogh i~Vedelsby~\cite{krogh1995} pokazali, że w~regresji błąd kwadratowy zespołu uśredniającego jest równy średniemu błędowi modeli pomniejszonemu o~rozrzut ich predykcji wokół średniej (\emph{ambiguity}). Brown, Wyatt i~Tiňo~\cite{brown2005} zwrócili uwagę, że różnorodność poprawia zespół tylko wtedy, gdy nie pogarsza nadmiernie modeli składowych.

\paragraph{Miary różnorodności.} Kuncheva i~Whitaker~\cite{kuncheva2003} porównali dziesięć miar różnorodności klasyfikatorów i~nie znaleźli jednoznacznego związku między ich wartościami a~dokładnością zespołu. Brown i~in.~\cite{brown2005survey} opisali losowanie obserwacji jako jedno z~podstawowych źródeł różnorodności modeli.

\paragraph{Luka badawcza.} Opisane miary dotyczą predykcji lub błędów nauczonych modeli. Właściwości samych próbek pojawiają się w~literaturze głównie pośrednio, jako rozmiar próbki. W~projekcie sprawdzamy, czy proste miary obliczone na próbkach przed uczeniem niosą informację o~jakości zespołu, przy wspólnym protokole dla klasyfikacji i~regresji.

% ====================================================================
\section{Metoda}
% ====================================================================

\subsection{Schemat badania}
\label{sec:schemat}

Dla każdego zbioru danych i~każdego z~50 podziałów na część treningową i~testową:
\begin{enumerate}
  \item dopasowujemy kodowanie cech kategorycznych wyłącznie na części treningowej;
  \item dla każdego z~15 wariantów losowania losujemy $M=100$ próbek indeksów części treningowej;
  \item obliczamy miary próbek \textbf{przed uczeniem drzew} i~bez dostępu do części testowej: miary różnorodności dla każdej z~$\binom{100}{2}=4950$ par próbek oraz miary reprezentatywności dla każdej ze 100 próbek;
  \item uczymy po jednym drzewie na każdej próbce i~łączymy ich predykcje;
  \item oceniamy pojedyncze drzewa, pary drzew i~cały zespół na części testowej.
\end{enumerate}
Liczba $M=100$ oznacza liczbę drzew (i~próbek) w~zespole, a~nie rozmiar próbki. Rozmiar próbki $m$ jest czynnikiem eksperymentalnym. W~wariancie odniesienia (S1, $r=1$, Sekcja~\ref{sec:warianty}) wynosi on $m=n$, czyli tyle, ile obserwacji liczy część treningowa.

\subsection{Model bazowy i~agregacja}
\label{sec:model}

W~eksperymencie głównym używamy lasu losowego w~domyślnej konfiguracji biblioteki \texttt{scikit-learn}~\cite{sklearn}. Drzewa rosną bez ograniczenia głębokości. W~klasyfikacji w~każdym węźle losowanych jest $\sqrt{d}$ cech, a~w~regresji brane są pod uwagę wszystkie cechy (\texttt{max\_features=1.0}), więc domyślny las losowy w~regresji jest baggingiem drzew. Wbudowane losowanie próbek wyłączyliśmy i~uczymy drzewa na próbkach generowanych przez nasz kod, z~zachowaniem powtórzeń obserwacji. Dzięki temu znamy każdą próbkę i~możemy obliczyć jej miary.

Dla klasyfikacji powtórzyliśmy cały eksperyment z~drzewami, które w~każdym węźle rozważają wszystkie cechy. Odpowiada to baggingowi w~pierwotnym sformułowaniu Breimana~\cite{breiman1996}. W~dalszej części tę konfigurację nazywamy \emph{baggingiem}, a~konfigurację główną \emph{lasem losowym}. W~baggingu jedynym źródłem różnic między drzewami jest losowanie obserwacji, co odpowiada pierwotnemu planowi projektu. Porównanie obu konfiguracji pokazuje, jak losowanie cech zmienia wpływ losowania próbek.

\textbf{Agregacja.} W~klasyfikacji zespół uśrednia prawdopodobieństwa klas przewidywane przez drzewa i~wybiera klasę o~największej średniej (głosowanie miękkie), tak jak las losowy w~\texttt{scikit-learn}. Wynik głosowania większościowego zapisujemy pomocniczo. W~regresji zespół zwraca średnią predykcji drzew.

\subsection{Warianty losowania próbek}
\label{sec:warianty}

Niech $n$ oznacza liczbę obserwacji części treningowej, a~$m$ rozmiar jednej próbki. Porównujemy 15 wariantów z~czterech rodzin (Tabela~\ref{tab:warianty}). Punktem odniesienia jest wariant S1 z~$r=1$, czyli losowanie ze zwracaniem $m=n$ obserwacji, stosowane domyślnie w~baggingu i~lesie losowym. W~dalszej części nazywamy go \emph{wariantem odniesienia}.

\begin{table}[!htbp]
\centering\small
\caption{Warianty losowania próbek.}
\label{tab:warianty}
\begin{tabularx}{\textwidth}{>{\raggedright\arraybackslash}p{0.2\textwidth}>{\raggedright\arraybackslash}p{0.3\textwidth}X}
\toprule
\textbf{Rodzina} & \textbf{Parametry} & \textbf{Co zmienia} \\
\midrule
S1: losowanie ze zwracaniem & $m=rn$, $r\in\{0.2,0.4,0.6,0.8,1\}$ & rozmiar próbki, a~przez to jednocześnie siłę drzew i~ich różnorodność \\
S2: losowanie bez zwracania & $m=sn$, $s\in\{0.18,0.33,0.45,0.55,0.63\}$ & to samo co S1 przy tej samej oczekiwanej liczbie różnych obserwacji ($s\approx1-e^{-r}$), ale bez powtórzeń \\
S3: wspólny rdzeń & bez zwracania, $m=0.5n$, $\gamma\in\{0,0.25,0.5,0.75\}$ & wyłącznie stopień pokrywania się próbek przy stałym rozmiarze \\
S4: losowanie warstwowe & ze zwracaniem, $m=n$, osobno w~każdej klasie (w~regresji w~10 przedziałach kwantylowych $y$) & usuwa losowe wahania proporcji klas lub rozkładu $y$ między próbkami \\
\bottomrule
\end{tabularx}
\end{table}

\textbf{S1 i~S2.} W~S1 każde drzewo losuje ze zwracaniem $m=rn$ obserwacji. Próbka zawiera wtedy około $1-e^{-r}$ różnych obserwacji części treningowej (dla $r=1$ około 63\%), a~pozostałe pozycje to powtórzenia. Wartości $s$ w~S2 dobraliśmy tak, aby próbka losowana bez zwracania zawierała tyle samo różnych obserwacji. Porównanie S1 z~S2 przy tym samym udziale różnych obserwacji pokazuje więc, czy na wynik wpływa samo powtarzanie obserwacji.

\textbf{S3: wspólny rdzeń.} Parametr $\gamma$ jest udziałem wspólnego rdzenia w~próbce. W~planie oznaczaliśmy go $\alpha$; zmieniliśmy oznaczenie, aby nie mylić go z~poziomem istotności. Najpierw, raz dla całego zespołu, losujemy bez zwracania $\gamma m$ obserwacji, które trafiają do próbki \emph{każdego} drzewa. Następnie każde drzewo osobno dolosowuje bez zwracania $(1-\gamma)m$ obserwacji spośród pozostałych $n-\gamma m$. Każda próbka ma więc ten sam rozmiar $m=0.5n$ i~rozpatrywana osobno jest zwykłą losową podpróbką. Zmienia się tylko stopień pokrywania się próbek różnych drzew. Przy $\gamma=0$ próbki są losowane niezależnie, a~przy $\gamma=0.75$ trzy czwarte każdej próbki są wspólne dla wszystkich drzew. Wariant S3 pozwala zatem oddzielić wpływ niezależności próbek od wpływu siły drzew.

\textbf{S4: losowanie warstwowe.} W~każdej klasie losujemy ze zwracaniem tyle obserwacji, ile liczy ona w~części treningowej. Każda próbka ma więc dokładnie takie proporcje klas jak część treningowa. W~regresji klasy zastępują przedziały kwantylowe $y$ wyznaczone na części treningowej.

Nie stosujemy nadpróbkowania klas ani wag obserwacji, bo zmieniałyby badany mechanizm losowania. W~danym podziale wszystkie warianty korzystają z~tego samego ziarna generatora liczb losowych, dzięki czemu porównania wariantów są sparowane.

\subsection{Miary próbek obliczane przed uczeniem}
\label{sec:miary}

Niech $S_i$ oznacza zbiór różnych obserwacji w~próbce drzewa $i$, a~$c_{ir}$ liczbę wystąpień obserwacji $r$ w~tej próbce. Miary par uśredniamy po 4950 parach drzew, a~miary pojedynczych próbek po 100 drzewach.

\textbf{Różnorodność par próbek} (odpowiednik niezależności drzew):
\begin{itemize}
  \item odległość Jaccarda $D_J(i,j)=1-\frac{|S_i\cap S_j|}{|S_i\cup S_j|}$, czyli część obserwacji występujących w~którejkolwiek z~dwóch próbek, która nie jest im wspólna;
  \item odległość krotności $D_C(i,j)=\frac12\sum_{r=1}^{n}\left|\frac{c_{ir}}{m}-\frac{c_{jr}}{m}\right|$, która uwzględnia także liczbę powtórzeń;
  \item różnica rozkładu zmiennej docelowej $D_Y(i,j)$. W~klasyfikacji jest to odległość całkowitej zmienności między proporcjami klas, $\frac12\sum_c|p_{ic}-p_{jc}|$. W~regresji jest to odległość Wassersteina między rozkładami $y$ w~obu próbkach, podzielona przez odchylenie standardowe $y$ w~części treningowej.
\end{itemize}

\textbf{Reprezentatywność pojedynczej próbki} (odpowiednik siły drzewa):
\begin{itemize}
  \item $U_i=|S_i|/n$, czyli udział różnych obserwacji części treningowej w~próbce;
  \item $R_{Y,i}$, czyli odległość rozkładu $y$ w~próbce od rozkładu w~całej części treningowej (te same odległości co w~$D_Y$);
  \item $R_{X,i}$, czyli średnia po cechach odległość Wassersteina między rozkładem cechy w~próbce i~w~części treningowej, liczona na cechach standaryzowanych parametrami części treningowej.
\end{itemize}
Mniejsza wartość $R_Y$ lub $R_X$ oznacza próbkę bardziej reprezentatywną. Wszystkie miary zależą wyłącznie od wylosowanych indeksów i~danych treningowych. Implementację wektorową (macierz krotności o~wymiarach $M\times n$) sprawdza test zgodności z~prostą implementacją liczoną para po parze (\texttt{tests/test\_diversity\_vectorized.py}).

\subsection{Wielkości mierzone po uczeniu}

Na części testowej mierzymy:
\begin{itemize}
  \item jakość każdego drzewa (macro-F1 albo RMSE i~$R^2$) oraz jej średnią i~odchylenie standardowe w~zespole, jako odpowiednik siły;
  \item różnorodność predykcji każdej pary drzew, jako odpowiednik niezależności. W~klasyfikacji jest to niezgodność~\cite{kuncheva2003}, czyli odsetek obserwacji, na których myli się dokładnie jedno z~dwóch drzew. W~regresji jest to $1-$korelacja błędów obu drzew;
  \item zysk z~agregacji, czyli różnicę między jakością zespołu a~średnią jakością drzew. W~regresji z~błędem kwadratowym zysk ten jest równy rozrzutowi predykcji drzew wokół ich średniej~\cite{krogh1995}. Zapisujemy obie wielkości, a~ich równość służy jako dodatkowy test poprawności kodu.
\end{itemize}

\subsection{Koncentracja miar próbek bootstrapowych}
\label{sec:koncentracja}

W~wariancie odniesienia miary próbek mają prawie zawsze te same wartości. Nie jest to błąd implementacji, lecz własność losowania bootstrapowego. Liczba wystąpień obserwacji w~próbce o~rozmiarze $m=rn$ losowanej ze zwracaniem ma rozkład $\mathrm{Bin}(m,1/n)\approx\mathrm{Poisson}(r)$, więc obserwacja trafia do próbki z~prawdopodobieństwem $p=1-(1-1/n)^{m}\to1-e^{-r}$. Ponieważ $U_i$ jest średnią $n$ prawie niezależnych wskaźników, $\mathbb{E}U_i\to1-e^{-r}$, a~odchylenie standardowe $U_i$ maleje jak $1/\sqrt{n}$ (dla $r=1$ wynosi około $0.31/\sqrt{n}$). To samo dotyczy pary próbek. Obserwacja trafia do obu próbek z~prawdopodobieństwem $p^2$, a~do co najmniej jednej z~prawdopodobieństwem $2p-p^2$, zatem
\[
D_J(i,j)\;\to\;1-\frac{p^2}{2p-p^2}=1-\frac{p}{2-p}\approx0.538\qquad(r=1),
\]
z~rozrzutem rzędu $1/\sqrt{n}$. Rysunek~\ref{fig:koncentracja} potwierdza to empirycznie. Przy kilkuset obserwacjach wszystkie pary próbek są prawie jednakowo odległe, a~przy kilku tysiącach różnice niemal znikają. Dlatego różnorodność próbek trzeba zmieniać celowo, wariantami S1--S4. Pytanie o~wpływ konkretnego losowania (eksperyment dodatkowy, Sekcja~\ref{sec:losowania}) dotyczy natomiast różnic, które z~natury są bardzo małe.

\begin{figure}[!htbp]
\centering
\includegraphics[width=\textwidth]{01_concentration.png}
\caption{Koncentracja miar próbek losowanych ze zwracaniem przy $m=n$. (a) Rozkład odległości Jaccarda dla 4950 par próbek przy rosnącym $n$; linia przerywana oznacza granicę $1-p/(2-p)$. (b) Odchylenie standardowe $D_J$ między parami w~jednym zespole oraz odchylenie standardowe średniego $D_J$ między podziałami maleją jak $1/\sqrt{n}$.}
\label{fig:koncentracja}
\end{figure}

\subsection{Eksperyment dodatkowy: wpływ konkretnego losowania}
\label{sec:losowania}

Sprawdzamy także, czy przy ustalonym wariancie na jakość zespołu wpływa to, które konkretnie próbki wylosowano. Dla każdego zbioru w~5 podziałach (pierwsze powtórzenie walidacji krzyżowej) uczymy 20 zespołów w~wariancie odniesienia, w~konfiguracji bagging. Zespoły różnią się wyłącznie ziarnem losowania próbek, a~ziarno drzew jest wspólne. Następnie w~każdym podziale obliczamy korelację Spearmana między miarą próbek a~jakością zespołu dla 20 losowań.

% ====================================================================
\section{Dane}
% ====================================================================

\subsection{Zbiory danych}

Badanie obejmuje 12 zbiorów klasyfikacyjnych i~12 regresyjnych z~repozytoriów UCI~\cite{uci} i~OpenML~\cite{openml} oraz z~biblioteki \texttt{scikit-learn} (Tabela~\ref{tab:zbiory}). Trzy zbiory pochodzą z~\texttt{scikit-learn}, dwa (Online Shoppers i~Combined Cycle Power Plant) pobieramy bezpośrednio z~UCI, a~pozostałe z~OpenML według ustalonego identyfikatora wersji. Zbiory różnią się liczbą obserwacji (od 208 do 20\,000), liczbą cech (od 4 do 81), liczbą klas (od 2 do 6) i~stopniem niezrównoważenia klas.

\begin{table}[!htbp]
\centering\scriptsize\setlength{\tabcolsep}{4pt}
\caption{Zbiory danych. $n_0$ oznacza liczbę obserwacji w~źródle, kolumny ,,braki'' i~,,duplikaty'' podają liczbę usuniętych wierszy, $n$ to liczba obserwacji użytych w~eksperymencie (zbiory większe niż 20\,000 obserwacji ograniczono losowo do 20\,000), a~$d$ to liczba cech przed kodowaniem one-hot i~po nim. Ostatnia kolumna zawiera liczebności klas po przygotowaniu danych.}
\label{tab:zbiory}
\begin{tabular}{llrrrrrl}
\toprule
\textbf{Zbiór} & \textbf{Źródło} & $n_0$ & \textbf{braki} & \textbf{duplikaty} & $n$ & $d$ & \textbf{klasy} \\
\midrule
\multicolumn{8}{l}{\emph{Klasyfikacja}}\\
Adult & OpenML & 48\,842 & 3620 & 163 & 20\,000 & 14 / 120 & 2 kl.: 15035/4965\\
Bank Marketing & OpenML & 45\,211 & -- & 16 & 20\,000 & 15 / 50 & 2 kl.: 17659/2341\\
Breast Cancer Wisconsin & sklearn & 569 & -- & -- & 569 & 30 / 30 & 2 kl.: 357/212\\
German Credit & OpenML & 1\,000 & -- & -- & 1\,000 & 20 / 61 & 2 kl.: 700/300\\
Glass & OpenML & 214 & -- & 1 & 213 & 9 / 9 & 6 kl.: 76/69/29/17/13/9\\
Ionosphere & OpenML & 351 & -- & 1 & 350 & 33 / 33 & 2 kl.: 225/125\\
Online Shoppers & UCI 468 & 12\,330 & -- & 125 & 12\,205 & 17 / 29 & 2 kl.: 10297/1908\\
Phoneme & OpenML & 5\,404 & -- & 9 & 5\,395 & 5 / 5 & 2 kl.: 3818/1577\\
Sonar & OpenML & 208 & -- & -- & 208 & 60 / 60 & 2 kl.: 111/97\\
Spambase & OpenML & 4\,601 & -- & 391 & 4\,210 & 57 / 57 & 2 kl.: 2531/1679\\
Vehicle & OpenML & 846 & -- & -- & 846 & 18 / 18 & 4 kl.: 218/217/212/199\\
Waveform-5000 & OpenML & 5\,000 & -- & -- & 5\,000 & 40 / 40 & 3 kl.: 1692/1655/1653\\
\midrule
\multicolumn{8}{l}{\emph{Regresja}}\\
Abalone & OpenML & 4\,177 & -- & -- & 4\,177 & 8 / 10 & \\
Airfoil Self-Noise & OpenML & 1\,503 & -- & -- & 1\,503 & 5 / 5 & \\
Auto MPG & OpenML & 398 & 6 & -- & 392 & 7 / 9 & \\
California Housing & sklearn & 20\,640 & -- & -- & 20\,000 & 8 / 8 & \\
Concrete & OpenML & 1\,030 & -- & 25 & 1\,005 & 8 / 8 & \\
Diabetes & sklearn & 442 & -- & -- & 442 & 10 / 10 & \\
Energy Efficiency (Y1) & OpenML & 768 & -- & -- & 768 & 8 / 8 & \\
Kin8nm & OpenML & 8\,192 & -- & -- & 8\,192 & 8 / 8 & \\
Combined Cycle Power Plant & UCI 294 & 9\,568 & -- & 41 & 9\,527 & 4 / 4 & \\
Superconductivity & OpenML & 21\,263 & -- & 66 & 20\,000 & 81 / 81 & \\
Wine Quality -- Red & OpenML & 1\,599 & -- & 240 & 1\,359 & 11 / 11 & \\
Yacht Hydrodynamics & OpenML & 308 & -- & -- & 308 & 6 / 6 & \\
\bottomrule
\end{tabular}
\end{table}


\subsection{Przygotowanie danych i~ochrona przed wyciekiem informacji}
\label{sec:prep}

Część operacji wykonujemy raz, na całym zbiorze, przed jakimkolwiek podziałem. Pozostałe wykonujemy w~każdym podziale osobno i~dopasowujemy wyłącznie na części treningowej. O~przydziale decyduje jedno kryterium. Przed podziałem wolno wykonać tylko operacje, które nie korzystają ze statystyk obliczonych na innych obserwacjach, czyli decyzje dotyczące pojedynczego wiersza albo całej kolumny, niezależne od tego, które obserwacje trafią do części testowej. Wszystko, co jest szacowane z~danych, dopasowujemy na części treningowej~\cite{kaufman2012}. Tabela~\ref{tab:prep} zestawia wszystkie kroki.

\begin{table}[!htbp]
\centering\small
\caption{Kolejność kroków przygotowania danych. Kroki 1--7 wykonujemy raz dla zbioru, a~kroki 8--11 osobno w~każdym z~50 podziałów.}
\label{tab:prep}
\begin{tabularx}{\textwidth}{r>{\raggedright\arraybackslash}p{0.42\textwidth}>{\raggedright\arraybackslash}X}
\toprule
& \textbf{Krok} & \textbf{Dlaczego nie powoduje wycieku} \\
\midrule
\multicolumn{3}{l}{\emph{Przed podziałem (cały zbiór)}}\\
1 & usunięcie kolumn identyfikujących obserwacje i~cech znanych dopiero po fakcie (\texttt{duration} w~Bank Marketing, nazwa samochodu w~Auto MPG, druga zmienna docelowa Y2 w~Energy Efficiency) & decyzja na podstawie dokumentacji zbioru, a~nie wartości w~danych \\
2 & usunięcie wierszy bez wartości zmiennej docelowej & dotyczy pojedynczego wiersza \\
3 & usunięcie wierszy z~brakiem w~dowolnej cesze (po zamianie znaczników takich jak ,,?'' na brak) & decyzja zależy tylko od samego wiersza; bez imputacji nie jest obliczana żadna statystyka \\
4 & usunięcie dokładnych duplikatów (te same cechy i~ta sama wartość zmiennej docelowej) & zapobiega trafieniu tej samej obserwacji jednocześnie do części treningowej i~testowej \\
5 & usunięcie klas liczących mniej niż 5 obserwacji & element definicji zadania; w~naszych zbiorach nie usunięto żadnej klasy \\
6 & losowe ograniczenie do 20\,000 obserwacji (w~klasyfikacji z~zachowaniem proporcji klas) & losowanie wierszy; wyznacza populację, z~której tworzymy podziały \\
7 & usunięcie kolumn stałych & kolumna stała w~całym zbiorze jest stała w~każdej części treningowej i~drzewo nigdy jej nie użyje; w~naszych zbiorach nie usunięto żadnej kolumny \\
\midrule
\multicolumn{3}{l}{\emph{W~każdym podziale (dopasowanie tylko na części treningowej)}}\\
8 & kodowanie one-hot cech kategorycznych & kategorie wyznaczone na części treningowej; kategoria nieobecna w~części treningowej jest kodowana jako wektor zer \\
9 & kodowanie klas liczbami & słownik klas z~części treningowej \\
10 & standaryzacja cech dla miary $R_X$, odchylenie standardowe $y$ dla $D_Y$ i~$R_Y$, przedziały kwantylowe $y$ dla S4 & parametry wyznaczone na części treningowej \\
11 & losowanie próbek, obliczanie miar próbek, uczenie drzew & tylko indeksy części treningowej \\
\bottomrule
\end{tabularx}
\end{table}

\textbf{Braki danych: usuwanie zamiast imputacji.} W~planie braki uzupełnialiśmy medianą (cechy liczbowe) lub osobną kategorią ,,brak'' (cechy kategoryczne). Imputacja dopasowana na części treningowej nie powoduje wycieku informacji, ale w~tym badaniu ma dwie wady. Po pierwsze, sztuczna kategoria i~wiele identycznych wartości mediany zmieniają rozkład cech, a~więc także miarę $R_X$. Po drugie, kategoria ,,brak'' może sama nieść informację o~klasie, której nie ma w~rzeczywistych danych. Dlatego ostatecznie usuwamy wiersze z~brakami. Braki występują tylko w~dwóch zbiorach (Tabela~\ref{tab:zbiory}): w~Adult w~7.4\% wierszy (cechy \texttt{workclass}, \texttt{occupation}, \texttt{native-country}) i~w~Auto MPG w~6 wierszach (\texttt{horsepower}). W~obu zbiorach jest to postępowanie standardowe; wersja zbioru Adult bez braków udostępniana w~repozytorium UCI liczy 45\,222 wiersze. Ponieważ Adult i~tak ograniczamy do 20\,000 obserwacji, rozmiar zbioru użytego w~eksperymencie się nie zmienia. Usuwanie wierszy zakłada, że braki występują całkowicie losowo (MCAR)~\cite{little2019}. W~Adult brak informacji o~zawodzie dotyczy częściej osób, które nie pracują, więc populacja po usunięciu wierszy jest nieco węższa od pierwotnej. Nie wpływa to na porównanie wariantów, bo wszystkie oceniamy na tych samych danych. Wartość \texttt{unknown} w~Bank Marketing jest udokumentowaną kategorią odpowiedzi, a~nie brakiem, więc tych wierszy nie usuwamy. Przetwarzanie danych w~podziale nie zawiera imputacji i~zgłasza błąd, gdyby brak dotarł do tego etapu.

\textbf{Duplikaty.} Kilka zbiorów zawiera wiersze identyczne łącznie ze zmienną docelową; najwięcej jest ich w~Spambase (391), Wine Quality (240) i~Online Shoppers (125). Przy losowym podziale kopie tej samej obserwacji mogą trafić jednocześnie do części treningowej i~testowej, co zawyża wynik na części testowej. Ponadto miary $U$ i~$D_J$ traktowałyby kopie jako różne obserwacje. Usuwamy więc duplikaty i~zostawiamy pierwsze wystąpienie. Wiersze o~tych samych cechach i~różnej wartości zmiennej docelowej zostają, bo opisują rzeczywisty szum (np.\ ten sam materiał w~Superconductivity zmierzony kilkakrotnie).

\textbf{Pozostałe zasady.} Cech liczbowych nie standaryzujemy, bo drzewa decyzyjne tego nie wymagają, a~zmienna docelowa w~regresji pozostaje w~oryginalnej skali. Nierównowagi klas nie korygujemy, ponieważ zmieniałoby to badane próbki. Uwzględniamy ją przez warstwowe podziały, metryki niezależne od klasy dominującej (macro-F1, zrównoważona dokładność, ROC-AUC) oraz wariant S4 i~miarę $D_Y$.

% ====================================================================
\section{Procedura eksperymentalna}
\label{sec:procedura}
% ====================================================================

\textbf{Walidacja.} Stosujemy pięciokrotną walidację krzyżową powtórzoną 10 razy, co daje 50 podziałów na zbiór; w~klasyfikacji podziały są warstwowe. Przy pięciu częściach zamiast dziesięciu część testowa w~najmniejszych zbiorach (Sonar, Glass) liczy ponad 40 obserwacji, a~dziesięć powtórzeń zmniejsza wpływ przypadkowego podziału. Wszystkie warianty oceniamy na tych samych podziałach i~z~tymi samymi ziarnami, więc porównania są sparowane. Łącznie nauczyliśmy $24\times50\times15=18\,000$ zespołów po 100 drzew w~eksperymencie głównym, $12\times50\times15=9000$ w~eksperymencie z~baggingiem oraz $24\times5\times20=2400$ w~eksperymencie dodatkowym.

\textbf{Pilotaż.} Przed eksperymentem głównym przeprowadziliśmy pilotaż na wszystkich 24 zbiorach z~jednym powtórzeniem walidacji (CV $5\times1$, te same 15 wariantów, las losowy i~bagging, 20 losowań w~eksperymencie dodatkowym). Jego wyniki przedstawiliśmy na punkcie kontrolnym 10.11.2026. Pilotaż posłużył do sprawdzenia poprawności kodu i~implementacji miar oraz do oszacowania czasu obliczeń (około 1 godziny procesora na powtórzenie). Jego wyników nie używamy w~analizie, ponieważ dane przygotowywaliśmy wtedy jeszcze z~imputacją braków i~bez usuwania duplikatów, a~eksperyment dodatkowy wykorzystywał las losowy. Wszystkie liczby w~raporcie pochodzą z~eksperymentu głównego.

\textbf{Hiperparametry} (Tabela~\ref{tab:hiper}) ustaliliśmy z~góry i~nie stroimy ich, ponieważ badamy wpływ losowania próbek, a~nie szukamy najlepszego modelu. Strojenie w~każdym podziale byłoby dodatkowym źródłem różnic między zespołami. Parametry losowania są czynnikiem eksperymentalnym, dlatego oceniamy wszystkie ich wartości.

\begin{table}[!htbp]
\centering\small
\setlength{\tabcolsep}{4pt}
\caption{Hiperparametry drzew (wartości domyślne biblioteki \texttt{scikit-learn}).}
\label{tab:hiper}
\begin{tabular}{lll}
\toprule
\textbf{Parametr} & \textbf{Las losowy (eksperyment główny)} & \textbf{Bagging (klasyfikacja)} \\
\midrule
liczba drzew $M$ & 100 & 100 \\
\texttt{max\_features} & $\sqrt{d}$ (klasyfikacja), wszystkie (regresja) & wszystkie \\
\texttt{max\_depth}, \texttt{min\_samples\_leaf} & \texttt{None}, 1 & \texttt{None}, 1 \\
kryterium podziału & \texttt{gini} / \texttt{squared\_error} & \texttt{gini} \\
agregacja & średnia prawdopodobieństw / średnia & średnia prawdopodobieństw \\
\bottomrule
\end{tabular}
\end{table}

\textbf{Metryki.} W~klasyfikacji metryką główną jest macro-F1. Zapisujemy też zrównoważoną dokładność, ROC-AUC (jedna klasa przeciw pozostałym, średnia z~równymi wagami~\cite{fawcett2006}) i~dokładność. W~regresji metryką główną jest RMSE; zapisujemy też MAE i~$R^2$. Testy statystyczne i~tabele w~raporcie dotyczą tylko metryk głównych.

% ====================================================================
\section{Analiza statystyczna}
\label{sec:stat}
% ====================================================================

\textbf{Jednostka analizy.} Wyniki z~50 podziałów nie są niezależne, bo części treningowe się pokrywają, więc testy na poziomie podziałów zaniżałyby wartości $p$ \cite{nadeau2003}. Jednostką w~testach jest zbiór danych. Wyniki z~50 podziałów sprowadzamy do jednej wartości na zbiór, średniej lub mediany~\cite{demsar2006}. Klasyfikację i~regresję analizujemy osobno, a~dodatkowo łącznie.

\textbf{Wynik względny.} Aby porównywać zbiory o~różnych skalach, wynik wariantu wyrażamy względem wariantu odniesienia w~tym samym zbiorze. W~klasyfikacji jest to różnica macro-F1 w~punktach procentowych, a~w~regresji względna poprawa RMSE, $(\mathrm{RMSE}_{S1,r=1}-\mathrm{RMSE})/\mathrm{RMSE}_{S1,r=1}$. Wartość dodatnia oznacza wariant lepszy od wariantu odniesienia.

\textbf{Testy.}
\begin{description}[leftmargin=1.2cm,style=sameline]
  \item[H1] W~każdym zespole wariantu odniesienia obliczamy korelację Spearmana po 4950 parach drzew (miary $D_J$, $D_C$, $D_Y$ a~różnorodność predykcji pary) i~po 100 drzewach (miary $U$, $R_Y$, $R_X$ a~jakość drzewa). Wynikiem zbioru jest mediana z~50 podziałów. Test Wilcoxona sprawdza, czy mediana po zbiorach różni się od zera.
  \item[H2] W~każdym zbiorze obliczamy korelację Spearmana między średnią wartością miary a~wynikiem zespołu po 15 wariantach oraz korelację cząstkową Spearmana, w~której kontrolujemy średnią jakość drzew. Wyniki zbiorów porównujemy testem Wilcoxona. Bezpośrednim testem drugiej części H2 jest porównanie wariantu S3 z~$\gamma=0$ i~z $\gamma=0.75$ testem Wilcoxona dla par.
  \item[H3] Test Friedmana z~poprawką Imana--Davenporta~\cite{demsar2006} sprawdza, czy 15 wariantów różni się jakością; podajemy też średnie rangi. Następnie każdy wariant porównujemy z~wariantem odniesienia testem Wilcoxona. Wykonujemy 14 porównań zamiast 105 porównań wszystkich par, zgodnie z~zaleceniem dla porównań z~metodą kontrolną~\cite{demsar2006}. Średnie wartości miar wariantów zestawiamy z~ich średnimi rangami.
\end{description}
Korektę Holma~\cite{holm1979} stosujemy osobno w~każdej rodzinie testów, np.\ dla 14 porównań z~wariantem odniesienia w~danym zadaniu. Poziom istotności wynosi $0.05$, a~wszystkie testy są dwustronne.

\textbf{Moc testów przy 12 zbiorach.} Przy 12 zbiorach najmniejsza osiągalna dwustronna wartość $p$ testu Wilcoxona wynosi $2/2^{12}\approx0.0005$, a~po korekcie Holma dla 14 testów najostrzejszy próg to $0.05/14\approx0.0036$. Aby wynik był istotny, wariant musi więc wygrać z~wariantem odniesienia w~prawie wszystkich zbiorach (suma rang porażek nie większa niż 4). Dlatego oprócz testów częstościowych:
\begin{itemize}
  \item analizujemy łącznie wszystkie 24 zbiory na wspólnym wyniku względnym, przy czym próg istotności jest wtedy znacznie łagodniejszy (Tabela~\ref{tab:moc});
  \item stosujemy bayesowski test znakowanych rang~\cite{benavoli2017} z~obszarem praktycznej równoważności (ROPE) równym $\pm1$\pp{} macro-F1 lub $\pm1\%$ RMSE. Implementacja jest własna, wzorowana na bibliotece \texttt{baycomp}~\cite{baycomp}, z~parametrem rozkładu a~priori $s=0.5$ jak w~tej bibliotece. Test podaje prawdopodobieństwa a~posteriori, że wariant jest gorszy od wariantu odniesienia, praktycznie mu równoważny lub od niego lepszy. Pozwala więc odróżnić brak różnicy od braku dowodów;
  \item podajemy wielkości efektów (medianę różnicy, liczbę wygranych, dwuseryjny współczynnik korelacji rang) oraz znaczenie praktyczne różnicy, czyli to, czy przekracza ona odchylenie standardowe wyniku wariantu odniesienia między 50 podziałami;
  \item porównujemy stabilność wyników, czyli stosunek odchyleń standardowych metryki głównej wariantu i~wariantu odniesienia między podziałami. Wariant stabilniejszy może być preferowany nawet przy nieco gorszej średniej.
\end{itemize}

\begin{table}[!htbp]
\centering\small
\caption{Moc testu Wilcoxona przy 12 i~24 zbiorach. Podano najmniejszą osiągalną dwustronną wartość $p$ oraz największą sumę rang $T$ zbiorów o~przeciwnym znaku różnicy, przy której wynik jest jeszcze istotny w~najostrzejszym kroku procedury Holma (14 testów, próg $0.05/14$) i~bez korekty (próg $0.05$).}
\label{tab:moc}
\begin{tabular}{rrrr}
\toprule
\textbf{zbiory} & \textbf{min.\ $p$} & \textbf{maks.\ $T$ (Holm, krok 1)} & \textbf{maks.\ $T$ (bez korekty)} \\
\midrule
12 & 0.00049 & 4 & 13 \\
24 & 1.2e-07 & 51 & 81 \\
\bottomrule
\end{tabular}
\end{table}


% ====================================================================
\section{Wyniki}
\label{sec:wyniki}
% ====================================================================

Wyniki przedstawiamy od najogólniejszego pytania, czyli który sposób losowania daje najlepszy zespół (H3). Następnie opisujemy mechanizm, czyli wpływ rozmiaru i~pokrywania się próbek na siłę i~niezależność drzew, a~potem hipotezy o~miarach próbek (H1, H2). Na końcu omawiamy eksperyment dodatkowy i~porównanie lasu losowego z~baggingiem. Średnie i~odchylenia standardowe metryki głównej z~50 podziałów dla każdego zbioru i~wariantu zawiera Dodatek~\ref{app:szczegoly}.

% --------------------------------------------------------------------
\subsection{Porównanie wariantów z~wariantem odniesienia (H3)}
\label{sec:wyniki-h3}

Test Friedmana odrzuca hipotezę o~równoważności 15 wariantów w~obu zadaniach (Tabela~\ref{tab:friedman}), więc sposób losowania wpływa na jakość zespołu. Średnie rangi (Rysunek~\ref{fig:rangi}) pokazują jednak, że wariant odniesienia należy do najlepszych. W~regresji ma najlepszą średnią rangę (4.3). W~klasyfikacji lepsze rangi mają tylko losowanie warstwowe S4 (3.9) i~losowanie bez zwracania z~$s=0.63$ (4.0), czyli warianty najbardziej do niego podobne.

\begin{table}[!htbp]
\centering\small
\caption{Test Friedmana dla 15 wariantów (las losowy) i~jego wersja Imana--Davenporta.}
\label{tab:friedman}
\begin{tabular}{lrrrrr}
\toprule
\textbf{Zadanie} & \textbf{zbiory} & $\chi^2_F$ & $p$ & $F_{ID}$ & $p$ \\
\midrule
Klasyfikacja & 12 & 75.5 & <0.001 & 8.99 & <0.001 \\
Regresja & 12 & 76.9 & <0.001 & 9.29 & <0.001 \\
\bottomrule
\end{tabular}
\end{table}


\begin{figure}[!htbp]
\centering
\includegraphics[width=\textwidth]{05_average_ranks.png}
\caption{Średnie rangi 15 wariantów w~teście Friedmana (12 zbiorów w~każdym zadaniu, las losowy). Czarną ramką oznaczono wariant odniesienia S1~$r=1$.}
\label{fig:rangi}
\end{figure}

Porównania z~wariantem odniesienia zawiera Tabela~\ref{tab:kontrola}.

Żaden wariant nie jest istotnie lepszy od wariantu odniesienia, ani w~klasyfikacji, ani w~regresji, ani w~analizie łącznej 24 zbiorów. Dwa warianty o~najlepszych rangach w~klasyfikacji (S4 i~S2~$s=0.63$) poprawiają macro-F1 o~medianę $0.09$ i~$0.04$\pp{}. Test bayesowski przypisuje obu prawdopodobieństwo $1.00$ praktycznej równoważności z~wariantem odniesienia, więc wynik wskazuje na brak różnicy praktycznej, a~nie tylko na brak istotności.

Istotnie gorsze są warianty z~dużym wspólnym rdzeniem: S3 z~$\gamma\ge0.25$ w~regresji i~z $\gamma=0.75$ w~klasyfikacji. Przegrywają one w~11 lub 12 z~12 zbiorów, a~w~regresji mediana wzrostu RMSE wynosi od 2.8 do 8.6\%.

Małe próbki (S1 z~$r\le0.4$, S2 z~$s\le0.33$) przegrywają z~wariantem odniesienia w~9--11 z~12 zbiorów, a~w~regresji mediana straty sięga 10\% RMSE. Przy 12 zbiorach wynik nie jest jednak istotny po korekcie Holma ($p_H$ od 0.08 do 0.16). Zgadza się to z~analizą mocy (Tabela~\ref{tab:moc}): wystarczy kilka zbiorów, w~których mniejsze próbki pomagają, aby test nie odrzucił hipotezy zerowej. W~analizie łącznej 24 zbiorów te same warianty są istotnie gorsze ($p_H\le0.014$), a~test bayesowski daje prawdopodobieństwo od 0.93 do 1.00, że są gorsze o~więcej niż szerokość ROPE.

Strata przy najmniejszych próbkach przekracza odchylenie standardowe wyniku wariantu odniesienia między podziałami w~5 zbiorach klasyfikacyjnych i~w~8 regresyjnych. Dla wariantów najbliższych wariantowi odniesienia (S1~$r=0.8$, S2~$s\ge0.55$, S4) różnica nie przekracza tego odchylenia w~żadnym zbiorze.

Mediana stosunku odchyleń standardowych $\sigma/\sigma_0$ wynosi od 0.99 do 1.09 dla wszystkich wariantów z~wyjątkiem S3~$\gamma=0.75$ w~regresji (1.20). Żaden wariant nie jest więc wyraźnie stabilniejszy od wariantu odniesienia, a~duży wspólny rdzeń, oprócz gorszej średniej, zwiększa zmienność wyniku między podziałami.

{\scriptsize\setlength{\tabcolsep}{4pt}
\begin{longtable}{lrrrrrrcrr}
\caption{Porównanie wariantów z~wariantem odniesienia S1~$r=1$ (las losowy). Ranga to średnia ranga w~teście Friedmana (1 oznacza najlepszy z~15 wariantów). Mediana to mediana wyniku względnego po zbiorach; wartość dodatnia oznacza wariant lepszy od wariantu odniesienia. W/P to liczba zbiorów, w~których wariant wygrał i~przegrał. $r_{rb}$ to dwuseryjny współczynnik korelacji rang, a~$p_H$ to wartość $p$ po korekcie Holma (pogrubiono $p_H<0.05$). Kolumna Bayes zawiera prawdopodobieństwa a~posteriori, że wariant jest gorszy, praktycznie równoważny (ROPE $\pm1$) i~lepszy. Kolumna ,,Istotne prakt.'' podaje liczbę zbiorów, w~których różnica na korzyść i~na niekorzyść wariantu przekracza odchylenie standardowe wyniku wariantu odniesienia między 50 podziałami. $\sigma/\sigma_0$ to mediana po zbiorach stosunku odchyleń standardowych metryki głównej wariantu i~wariantu odniesienia między podziałami; wartość poniżej 1 oznacza wariant stabilniejszy.}
\label{tab:kontrola}\\
\toprule
\textbf{Wariant} & \textbf{Ranga} & \textbf{Mediana} & \textbf{W/P} & $r_{rb}$ & $p$ & $p_H$ & \textbf{Bayes} & \textbf{Istotne prakt.} & $\sigma/\sigma_0$ \\
\endfirsthead
\toprule
\textbf{Wariant} & \textbf{Ranga} & \textbf{Mediana} & \textbf{W/P} & $r_{rb}$ & $p$ & $p_H$ & \textbf{Bayes} & \textbf{Istotne prakt.} & $\sigma/\sigma_0$ \\
\endhead
\midrule
\multicolumn{10}{l}{\emph{Klasyfikacja} -- $\Delta$ macro-F1 [p.p.]}\\
S1~$r=0.2$ & 12.2 & -1.62 & 1/11 & -0.85 & 0.007 & 0.089 & 0.88 / 0.12 / 0.00 & 1/5 & 1.07 \\
S1~$r=0.4$ & 10.6 & -0.58 & 1/11 & -0.77 & 0.016 & 0.161 & 0.23 / 0.77 / 0.00 & 1/1 & 1.05 \\
S1~$r=0.6$ & 8.0 & -0.28 & 3/9 & -0.67 & 0.042 & 0.382 & 0.00 / 1.00 / 0.00 & 1/0 & 1.04 \\
S1~$r=0.8$ & 5.8 & -0.13 & 5/7 & -0.08 & 0.850 & 1.000 & 0.00 / 1.00 / 0.00 & 0/0 & 1.04 \\
S2~$s=0.18$ & 13.0 & -1.58 & 1/11 & -0.85 & 0.007 & 0.089 & 0.90 / 0.10 / 0.00 & 1/5 & 1.01 \\
S2~$s=0.33$ & 11.4 & -0.75 & 1/11 & -0.82 & 0.009 & 0.102 & 0.27 / 0.73 / 0.00 & 1/1 & 1.01 \\
S2~$s=0.45$ & 6.2 & -0.13 & 5/7 & -0.38 & 0.266 & 1.000 & 0.00 / 1.00 / 0.00 & 1/0 & 1.01 \\
S2~$s=0.55$ & 5.3 & -0.03 & 6/6 & -0.18 & 0.622 & 1.000 & 0.00 / 1.00 / 0.00 & 0/0 & 0.99 \\
S2~$s=0.63$ & 4.0 & +0.04 & 8/4 & +0.36 & 0.301 & 1.000 & 0.00 / 1.00 / 0.00 & 0/0 & 1.03 \\
S3~$\gamma=0$ & 7.8 & -0.23 & 3/9 & -0.62 & 0.064 & 0.512 & 0.00 / 1.00 / 0.00 & 0/0 & 1.06 \\
S3~$\gamma=0.25$ & 7.1 & -0.12 & 4/8 & -0.38 & 0.266 & 1.000 & 0.00 / 1.00 / 0.00 & 0/0 & 1.08 \\
S3~$\gamma=0.5$ & 8.7 & -0.19 & 3/9 & -0.62 & 0.064 & 0.512 & 0.00 / 1.00 / 0.00 & 0/1 & 1.04 \\
S3~$\gamma=0.75$ & 10.8 & \textbf{-0.48} & 1/11 & -0.97 & <0.001 & \textbf{0.014} & 0.18 / 0.82 / 0.00 & 0/1 & 1.09 \\
S4 & 3.9 & +0.09 & 8/4 & +0.62 & 0.064 & 0.512 & 0.00 / 1.00 / 0.00 & 0/0 & 1.03 \\
\midrule
\multicolumn{10}{l}{\emph{Regresja} -- poprawa RMSE [\%]}\\
S1~$r=0.2$ & 11.3 & -10.16 & 3/9 & -0.85 & 0.007 & 0.075 & 1.00 / 0.00 / 0.00 & 0/8 & 1.01 \\
S1~$r=0.4$ & 8.7 & -4.40 & 3/9 & -0.79 & 0.012 & 0.110 & 0.99 / 0.01 / 0.00 & 0/7 & 0.99 \\
S1~$r=0.6$ & 6.3 & -1.82 & 4/8 & -0.69 & 0.034 & 0.205 & 0.91 / 0.08 / 0.01 & 0/2 & 0.99 \\
S1~$r=0.8$ & 4.6 & -0.44 & 5/7 & -0.46 & 0.176 & 0.353 & 0.11 / 0.89 / 0.00 & 0/0 & 0.99 \\
S2~$s=0.18$ & 11.3 & -10.20 & 3/9 & -0.85 & 0.007 & 0.075 & 1.00 / 0.00 / 0.00 & 0/8 & 1.01 \\
S2~$s=0.33$ & 9.0 & -4.45 & 3/9 & -0.74 & 0.021 & 0.168 & 0.99 / 0.01 / 0.00 & 0/7 & 1.00 \\
S2~$s=0.45$ & 6.7 & -1.99 & 4/8 & -0.72 & 0.027 & 0.188 & 0.90 / 0.10 / 0.00 & 0/2 & 1.00 \\
S2~$s=0.55$ & 5.9 & -0.75 & 4/8 & -0.67 & 0.042 & 0.205 & 0.36 / 0.64 / 0.00 & 0/0 & 1.01 \\
S2~$s=0.63$ & 5.9 & -0.60 & 1/11 & -0.69 & 0.034 & 0.205 & 0.04 / 0.96 / 0.00 & 0/0 & 1.01 \\
S3~$\gamma=0$ & 6.2 & -1.07 & 4/8 & -0.67 & 0.042 & 0.205 & 0.57 / 0.43 / 0.00 & 0/2 & 1.00 \\
S3~$\gamma=0.25$ & 9.8 & \textbf{-2.77} & 1/11 & -0.97 & <0.001 & \textbf{0.012} & 0.98 / 0.02 / 0.00 & 0/2 & 1.02 \\
S3~$\gamma=0.5$ & 11.7 & \textbf{-4.54} & 0/12 & -1.00 & <0.001 & \textbf{0.007} & 1.00 / 0.00 / 0.00 & 0/7 & 1.06 \\
S3~$\gamma=0.75$ & 13.8 & \textbf{-8.64} & 0/12 & -1.00 & <0.001 & \textbf{0.007} & 1.00 / 0.00 / 0.00 & 0/7 & 1.20 \\
S4 & 4.6 & -0.03 & 4/8 & -0.23 & 0.519 & 0.519 & 0.00 / 1.00 / 0.00 & 0/0 & 1.01 \\
\midrule
\multicolumn{10}{l}{\emph{Łącznie (24 zbiory)} -- wynik względny [p.p. / \%]}\\
S1~$r=0.2$ & -- & \textbf{-3.74} & 4/20 & -0.79 & <0.001 & \textbf{0.003} & 1.00 / 0.00 / 0.00 & -- & -- \\
S1~$r=0.4$ & -- & \textbf{-1.05} & 4/20 & -0.70 & 0.002 & \textbf{0.014} & 0.93 / 0.07 / 0.00 & -- & -- \\
S1~$r=0.6$ & -- & -0.41 & 7/17 & -0.57 & 0.013 & 0.088 & 0.28 / 0.72 / 0.00 & -- & -- \\
S1~$r=0.8$ & -- & -0.17 & 10/14 & -0.25 & 0.290 & 0.579 & 0.00 / 1.00 / 0.00 & -- & -- \\
S2~$s=0.18$ & -- & \textbf{-3.85} & 4/20 & -0.79 & <0.001 & \textbf{0.003} & 1.00 / 0.00 / 0.00 & -- & -- \\
S2~$s=0.33$ & -- & \textbf{-1.18} & 4/20 & -0.71 & 0.001 & \textbf{0.013} & 0.94 / 0.06 / 0.00 & -- & -- \\
S2~$s=0.45$ & -- & -0.36 & 9/15 & -0.50 & 0.031 & 0.157 & 0.28 / 0.72 / 0.00 & -- & -- \\
S2~$s=0.55$ & -- & -0.11 & 10/14 & -0.47 & 0.042 & 0.170 & 0.02 / 0.98 / 0.00 & -- & -- \\
S2~$s=0.63$ & -- & -0.27 & 9/15 & -0.43 & 0.069 & 0.207 & 0.00 / 1.00 / 0.00 & -- & -- \\
S3~$\gamma=0$ & -- & -0.27 & 7/17 & -0.55 & 0.016 & 0.098 & 0.05 / 0.95 / 0.00 & -- & -- \\
S3~$\gamma=0.25$ & -- & \textbf{-0.52} & 5/19 & -0.78 & <0.001 & \textbf{0.004} & 0.56 / 0.44 / 0.00 & -- & -- \\
S3~$\gamma=0.5$ & -- & \textbf{-1.47} & 3/21 & -0.90 & <0.001 & \textbf{<0.001} & 0.92 / 0.08 / 0.00 & -- & -- \\
S3~$\gamma=0.75$ & -- & \textbf{-3.81} & 1/23 & -0.99 & <0.001 & \textbf{<0.001} & 1.00 / 0.00 / 0.00 & -- & -- \\
S4 & -- & +0.02 & 12/12 & +0.21 & 0.390 & 0.579 & 0.00 / 1.00 / 0.00 & -- & -- \\
\bottomrule
\end{longtable}
}


Rysunek~\ref{fig:zbiory} pokazuje wynik każdego wariantu w~każdym zbiorze. Kierunek wpływu rozmiaru próbki zależy od zbioru. Mniejsze próbki poprawiają wynik w~zbiorach z~dużym szumem w~zmiennej docelowej: Adult (do $+1.5$\pp{} macro-F1), Abalone ($+2.4\%$), Diabetes ($+3.2\%$) i~Wine Quality ($+0.6\%$). Znacznie pogarszają go natomiast w~zbiorach regresyjnych o~niemal deterministycznej zależności i~niewielu cechach. W~zbiorach Airfoil, Energy Efficiency, Yacht i~Concrete RMSE przy $r=0.2$ rośnie o~25--48\%. W~klasyfikacji największą stratę ma Glass ($-16.7$\pp), najmniejszy zbiór, w~którym jest sześć klas, a~najrzadsza z~nich ma w~części treningowej około 7 obserwacji. Próbka z~$r=0.2$ często w~ogóle jej nie zawiera. W~tym samym zbiorze losowanie warstwowe daje największą poprawę ($+1.3$\pp), bo zapewnia obecność każdej klasy w~każdej próbce.

\begin{figure}[!htbp]
\centering
\includegraphics[width=\textwidth]{08_per_dataset.png}
\caption{Wynik każdego wariantu względem wariantu odniesienia w~każdym zbiorze (średnia z~50 podziałów, las losowy). Lewy panel: różnica macro-F1 w~punktach procentowych. Prawy panel: względna poprawa RMSE w~procentach. Kolor niebieski oznacza wariant lepszy od wariantu odniesienia.}
\label{fig:zbiory}
\end{figure}

% --------------------------------------------------------------------
\subsection{Rozmiar próbki i~pokrywanie się próbek}
\label{sec:wyniki-rozmiar}

Rysunek~\ref{fig:rozmiar} pokazuje dwa czynniki, które warianty zmieniają celowo. Lewy panel przedstawia wynik S1 i~S2 w~funkcji udziału różnych obserwacji $U$. Krzywe obu rodzin prawie się pokrywają, więc przy tej samej liczbie różnych obserwacji nie ma znaczenia, czy próbka zawiera powtórzenia. Potwierdzają to mediany w~analizie łącznej (Tabela~\ref{tab:kontrola}). S1~$r=0.2$ i~S2~$s=0.18$ tracą odpowiednio 3.74 i~3.85 jednostki, a~S1~$r=0.4$ i~S2~$s=0.33$ odpowiednio 1.05 i~1.18 jednostki. Mediana wyniku rośnie monotonicznie z~$U$, choć, jak pokazano wyżej, nie w~każdym zbiorze.

\begin{figure}[!htbp]
\centering
\includegraphics[width=\textwidth]{02_size_and_core.png}
\caption{Lewy panel: wynik względem wariantu odniesienia w~funkcji udziału różnych obserwacji $U$ dla S1 (losowanie ze zwracaniem) i~S2 (bez zwracania). Prawy panel: wynik S3 w~funkcji udziału wspólnego rdzenia $\gamma$ przy stałym rozmiarze $m=0.5n$. Cienkie linie odpowiadają poszczególnym zbiorom, grube medianie po zbiorach.}
\label{fig:rozmiar}
\end{figure}

W~prawym panelu Rysunku~\ref{fig:rozmiar} rozmiar próbki jest stały ($m=0.5n$), a~zmienia się tylko stopień pokrywania się próbek różnych drzew. W~regresji jakość zespołu maleje monotonicznie z~$\gamma$ w~każdym zbiorze. W~klasyfikacji wyraźny spadek występuje dopiero przy $\gamma=0.75$. Tabela~\ref{tab:s3} zawiera bezpośredni test drugiej części H2. Próbki niezależne ($\gamma=0$) dają lepszy zespół niż próbki w~75\% wspólne w~9 z~12 zbiorów klasyfikacyjnych ($p=0.007$) i~we wszystkich 12 regresyjnych ($p<0.001$, mediana $+7.4\%$ RMSE). W~klasyfikacji różnica jest jednak mała: według testu bayesowskiego z~prawdopodobieństwem 0.94 mieści się w~ROPE, czyli jest mniejsza niż 1\pp{} macro-F1. Sekcja~\ref{sec:rf-bagging} pokazuje, że wynika to z~losowania cech w~lesie losowym.

\begin{table}[!htbp]
\centering\small
\caption{Bezpośredni test drugiej części H2. Wariant S3 z~$\gamma=0$ (próbki niezależne) porównano z~$\gamma=0.75$ (75\% każdej próbki wspólne dla wszystkich drzew) przy tym samym rozmiarze $m=0.5n$. Dodatnia różnica oznacza, że mniejsze pokrywanie się próbek daje lepszy zespół. Kolumna Bayes zawiera kolejno prawdopodobieństwa a~posteriori, że lepsze jest $\gamma=0.75$, że różnica mieści się w~ROPE i~że lepsze jest $\gamma=0$.}
\label{tab:s3}
\begin{tabular}{lrrrrc}
\toprule
\textbf{Zadanie} & \textbf{zbiory} & \textbf{mediana różnicy} & \textbf{wygrane $\gamma=0$} & $p$ & \textbf{Bayes} \\
\midrule
Klasyfikacja & 12 & +0.31 & 9 & 0.007 & 0.00 / 0.94 / 0.06 \\
Regresja & 12 & +7.38 & 12 & <0.001 & 0.00 / 0.00 / 1.00 \\
Łącznie (24 zbiory) & 24 & +4.20 & 21 & <0.001 & 0.00 / 0.00 / 1.00 \\
\bottomrule
\end{tabular}
\end{table}


% --------------------------------------------------------------------
\subsection{Siła i~niezależność drzew}
\label{sec:wyniki-sila}

Rysunek~\ref{fig:sila} pokazuje położenie zespołów na płaszczyźnie wyznaczonej przez siłę drzew i~różnorodność ich predykcji. Obie wielkości standaryzowaliśmy w~obrębie zbioru, a~jeden punkt odpowiada jednemu wariantowi w~jednym zbiorze. Warianty S1 i~S2 leżą na jednej malejącej krzywej: mniejsza próbka daje drzewa bardziej różne, ale słabsze. Warianty S3 układają się w~poziomy pas, ponieważ przy prawie stałej sile drzew zmienia się tylko ich różnorodność, a~zespoły o~mniejszej różnorodności (większe $\gamma$) są gorsze. Wariant S3 rozdziela więc, zgodnie z~założeniem, niezależność drzew od ich siły.

\begin{figure}[!htbp]
\centering
\includegraphics[width=\textwidth]{03_strength_vs_diversity.png}
\caption{Średnia jakość drzew a~różnorodność ich predykcji (wartości standaryzowane w~obrębie zbioru), las losowy. Każdy punkt odpowiada jednemu wariantowi w~jednym zbiorze, a~kolor wypełnienia wynikowi zespołu względem wariantu odniesienia.}
\label{fig:sila}
\end{figure}

% --------------------------------------------------------------------
\subsection{H1: miary próbek a~różnice między drzewami jednego zespołu}
\label{sec:wyniki-h1}

Tabela~\ref{tab:h1} i~Rysunek~\ref{fig:h1} zawierają korelacje między miarami próbek a~własnościami drzew w~obrębie jednego zespołu wariantu odniesienia.

Pary drzew uczonych na bardziej różnych próbkach ($D_J$, $D_C$) dają bardziej różne predykcje. Korelacja jest dodatnia w~10--12 z~12 zbiorów i~istotna po korekcie Holma w~obu zadaniach i~w~obu konfiguracjach zespołu. Jest jednak bardzo słaba. Mediana $\rho$ wynosi 0.02--0.03, więc miara wyjaśnia rzędu 0.1\% zmienności różnorodności par. Różnica rozkładu zmiennej docelowej $D_Y$ wiąże się z~różnorodnością predykcji jeszcze słabiej i~istotnie tylko w~regresji ($\tilde\rho=0.007$).

Żadna z~miar reprezentatywności ($U$, $R_Y$, $R_X$) nie jest istotnie związana z~jakością pojedynczego drzewa. Dla $U$ korelacja jest dodatnia w~9 lub 10 zbiorach ($\tilde\rho\approx0.02$), ale nie jest istotna po korekcie Holma. Dla $R_Y$ i~$R_X$ znak korelacji zmienia się między zbiorami.

\begin{table}[!htbp]
\centering\small
\caption{H1: korelacje Spearmana wewnątrz zespołu w~wariancie odniesienia S1~$r=1$. Podano medianę po zbiorach (z~median po 50 podziałach), liczbę zbiorów z~dodatnią korelacją oraz $p_H$, czyli wartość $p$ testu Wilcoxona po zbiorach z~korektą Holma (6 testów w~zadaniu). W~regresji las losowy jest baggingiem, więc kolumny dla baggingu pominięto.}
\label{tab:h1}
\begin{tabular}{lrrrrrr}
\toprule
& \multicolumn{3}{c}{\textbf{Las losowy}} & \multicolumn{3}{c}{\textbf{Bagging}} \\
\cmidrule(lr){2-4}\cmidrule(lr){5-7}
\textbf{Korelacja} & $\tilde\rho$ & $\rho>0$ & $p_H$ & $\tilde\rho$ & $\rho>0$ & $p_H$ \\
\midrule
\multicolumn{7}{l}{\emph{Klasyfikacja}}\\
pary: $D_J$ ~ różn.\ predykcji & +0.018 & 10/12 & \textbf{0.015} & +0.033 & 12/12 & \textbf{0.003} \\
pary: $D_C$ ~ różn.\ predykcji & +0.018 & 10/12 & \textbf{0.015} & +0.031 & 12/12 & \textbf{0.003} \\
pary: $D_Y$ ~ różn.\ predykcji & +0.003 & 9/12 & 0.904 & +0.006 & 7/12 & 0.529 \\
drzewa: $U$ ~ jakość drzewa & +0.017 & 9/12 & 0.107 & +0.017 & 9/12 & 0.369 \\
drzewa: $R_Y$ ~ jakość drzewa & +0.004 & 8/12 & 0.904 & -0.011 & 4/12 & 0.529 \\
drzewa: $R_X$ ~ jakość drzewa & -0.016 & 4/12 & 0.904 & -0.004 & 5/12 & 0.529 \\
\midrule
\multicolumn{7}{l}{\emph{Regresja}}\\
pary: $D_J$ ~ różn.\ predykcji & +0.033 & 12/12 & \textbf{0.003} & \multicolumn{3}{c}{= las losowy} \\
pary: $D_C$ ~ różn.\ predykcji & +0.031 & 12/12 & \textbf{0.003} & \multicolumn{3}{c}{= las losowy} \\
pary: $D_Y$ ~ różn.\ predykcji & +0.007 & 10/12 & \textbf{0.037} & \multicolumn{3}{c}{= las losowy} \\
drzewa: $U$ ~ jakość drzewa & +0.023 & 10/12 & 0.127 & \multicolumn{3}{c}{= las losowy} \\
drzewa: $R_Y$ ~ jakość drzewa & -0.003 & 6/12 & 0.761 & \multicolumn{3}{c}{= las losowy} \\
drzewa: $R_X$ ~ jakość drzewa & -0.002 & 5/12 & 0.910 & \multicolumn{3}{c}{= las losowy} \\
\bottomrule
\end{tabular}
\end{table}


\begin{figure}[!htbp]
\centering
\includegraphics[width=\textwidth]{04_h1_within_ensemble.png}
\caption{H1 dla lasu losowego w~wariancie odniesienia: mediana (po 50 podziałach) korelacji Spearmana w~każdym zbiorze. Kolor niebieski: miary par próbek a~różnorodność predykcji par drzew. Kolor pomarańczowy: miary pojedynczej próbki a~jakość drzewa.}
\label{fig:h1}
\end{figure}

Przyczynę słabości tych związków widać na Rysunku~\ref{fig:pary}. W~jednym zespole odległość Jaccarda wszystkich 4950 par próbek mieści się w~wąskim przedziale od około 0.46 do 0.62, zgodnie z~koncentracją opisaną w~Sekcji~\ref{sec:koncentracja}. Różnorodność predykcji tych samych par zmienia się natomiast kilkukrotnie. Różnice między drzewami zespołu wynikają więc przede wszystkim z~tego, \emph{które} obserwacje (a~w~lesie losowym także które cechy) trafiły do próbki, a~nie z~tego, \emph{ile} obserwacji dwie próbki mają wspólnych.

\begin{figure}[!htbp]
\centering
\includegraphics[width=\textwidth]{09_pair_example.png}
\caption{Przykład H1 w~jednym podziale: odległość Jaccarda i~różnorodność predykcji dla 4950 par drzew w~zbiorach Breast Cancer (klasyfikacja, niezgodność) i~Diabetes (regresja, $1-$korelacja błędów).}
\label{fig:pary}
\end{figure}

% --------------------------------------------------------------------
\subsection{H2: miary próbek a~jakość zespołu}
\label{sec:wyniki-h2}

Tabela~\ref{tab:h2} zawiera korelacje między średnią wartością miary a~jakością zespołu po 15 wariantach, liczone w~każdym zbiorze osobno. Rysunek~\ref{fig:h2} pokazuje te korelacje dla poszczególnych zbiorów.

\textbf{Korelacja prosta} ma w~większości zbiorów ten sam znak. Wyższe $U$ idzie w~parze z~lepszym zespołem, a~większa różnorodność par ($D_J$, $D_C$, $D_Y$) i~gorsza reprezentatywność ($R_Y$, $R_X$) z~gorszym. Wszystkie te miary zmieniają się bowiem razem z~rozmiarem próbki. Kierunek nie jest jednak wspólny dla wszystkich zbiorów. W~Adult, Abalone, Diabetes i~Wine Quality znaki są odwrócone, a~są to te same zbiory, w~których mniejsze próbki poprawiają wynik (Rysunek~\ref{fig:zbiory}). Dlatego żadna korelacja prosta nie jest istotna po korekcie Holma, choć np.\ dla $U$ znak jest ten sam w~9 lub 10 z~12 zbiorów.

\textbf{Korelacja cząstkowa}, czyli korelacja po usunięciu związku ze średnią jakością drzew, jest dodatnia prawie we wszystkich zbiorach (dolny wiersz Rysunku~\ref{fig:h2}). W~regresji wszystkie miary różnorodności par są dodatnio związane z~jakością zespołu we wszystkich 12 zbiorach ($\tilde\rho$ od 0.58 do 0.69, $p_H=0.003$). Najsilniejszy związek ma różnorodność predykcji ($\tilde\rho=0.77$). Przy tej samej sile drzew bardziej różne próbki dają więc lepszy zespół, zgodnie z~oszacowaniem Breimana. Dodatnie korelacje cząstkowe $R_Y$ i~$R_X$ mają tę samą przyczynę. Przy ustalonej sile drzew mniej reprezentatywna próbka jest w~praktyce próbką mniejszą, a~więc bardziej różną od pozostałych. W~lesie losowym w~klasyfikacji korelacje cząstkowe są dodatnie w~9 lub 10 z~12 zbiorów, ale słabsze i~nieistotne po korekcie ($p_H\ge0.19$). W~baggingu (Tabela~\ref{tab:h2bag}) ten sam związek jest istotny i~występuje we wszystkich 12 zbiorach ($\tilde\rho$ od 0.54 do 0.62 dla miar par, $p_H=0.003$).

\begin{table}[!htbp]
\centering\small
\caption{H2: korelacja Spearmana między średnią wartością miary a~wynikiem zespołu po 15 wariantach, liczona w~każdym zbiorze osobno. Podano medianę po zbiorach, liczbę zbiorów z~dodatnią korelacją i~wartość $p_H$ testu Wilcoxona z~korektą Holma (7 miar). Korelacja cząstkowa uwzględnia średnią jakość drzew.}
\label{tab:h2}
\begin{tabular}{lrrrrrr}
\toprule
& \multicolumn{3}{c}{\textbf{korelacja}} & \multicolumn{3}{c}{\textbf{cząstkowa}} \\
\cmidrule(lr){2-4}\cmidrule(lr){5-7}
\textbf{Miara} & $\tilde\rho$ & $\rho>0$ & $p_H$ & $\tilde\rho$ & $\rho>0$ & $p_H$ \\
\midrule
\multicolumn{7}{l}{\emph{Klasyfikacja}}\\
$D_J$ & -0.56 & 2/12 & 0.196 & +0.28 & 9/12 & 0.205 \\
$D_C$ & -0.45 & 2/12 & 0.196 & +0.22 & 10/12 & 0.212 \\
$D_Y$ & -0.60 & 2/12 & 0.196 & +0.27 & 9/12 & 0.330 \\
$U$ & +0.83 & 10/12 & 0.188 & +0.22 & 9/12 & 0.212 \\
$R_Y$ & -0.75 & 2/12 & 0.196 & +0.08 & 7/12 & 1.000 \\
$R_X$ & -0.69 & 2/12 & 0.196 & +0.04 & 7/12 & 1.000 \\
różnorodność predykcji & -0.64 & 2/12 & 0.196 & +0.41 & 10/12 & 0.188 \\
\midrule
\multicolumn{7}{l}{\emph{Regresja}}\\
$D_J$ & -0.42 & 4/12 & 1.000 & +0.59 & 12/12 & \textbf{0.003} \\
$D_C$ & -0.27 & 4/12 & 1.000 & +0.69 & 12/12 & \textbf{0.003} \\
$D_Y$ & -0.46 & 4/12 & 1.000 & +0.58 & 12/12 & \textbf{0.003} \\
$U$ & +0.83 & 9/12 & 0.113 & +0.18 & 10/12 & 0.233 \\
$R_Y$ & -0.60 & 3/12 & 1.000 & +0.44 & 12/12 & \textbf{0.003} \\
$R_X$ & -0.65 & 3/12 & 1.000 & +0.29 & 12/12 & \textbf{0.003} \\
różnorodność predykcji & -0.28 & 5/12 & 1.000 & +0.77 & 12/12 & \textbf{0.003} \\
\bottomrule
\end{tabular}
\end{table}


\begin{table}[!htbp]
\centering\small
\caption{H2: korelacja Spearmana między średnią wartością miary a~wynikiem zespołu po 15 wariantach, liczona w~każdym zbiorze osobno (bagging, klasyfikacja). Podano medianę po zbiorach, liczbę zbiorów z~dodatnią korelacją i~wartość $p_H$ testu Wilcoxona z~korektą Holma (7 miar). Korelacja cząstkowa uwzględnia średnią jakość drzew.}
\label{tab:h2bag}
\begin{tabular}{lrrrrrr}
\toprule
& \multicolumn{3}{c}{\textbf{korelacja}} & \multicolumn{3}{c}{\textbf{cząstkowa}} \\
\cmidrule(lr){2-4}\cmidrule(lr){5-7}
\textbf{Miara} & $\tilde\rho$ & $\rho>0$ & $p_H$ & $\tilde\rho$ & $\rho>0$ & $p_H$ \\
\midrule
\multicolumn{7}{l}{\emph{Klasyfikacja}}\\
$D_J$ & -0.16 & 5/12 & 1.000 & +0.54 & 12/12 & \textbf{0.003} \\
$D_C$ & -0.09 & 5/12 & 1.000 & +0.60 & 12/12 & \textbf{0.003} \\
$D_Y$ & -0.24 & 5/12 & 1.000 & +0.54 & 12/12 & \textbf{0.003} \\
$U$ & +0.50 & 8/12 & 1.000 & +0.12 & 8/12 & 0.339 \\
$R_Y$ & -0.36 & 4/12 & 1.000 & +0.26 & 9/12 & 0.068 \\
$R_X$ & -0.37 & 5/12 & 1.000 & +0.41 & 10/12 & \textbf{0.028} \\
różnorodność predykcji & -0.16 & 5/12 & 1.000 & +0.62 & 12/12 & \textbf{0.003} \\
\bottomrule
\end{tabular}
\end{table}


\begin{figure}[!htbp]
\centering
\includegraphics[width=\textwidth]{06_h2_correlations.png}
\caption{H2 w~poszczególnych zbiorach (las losowy). Górny wiersz: korelacja Spearmana między średnią wartością miary a~jakością zespołu po 15 wariantach. Dolny wiersz: korelacja cząstkowa przy kontroli średniej jakości drzew.}
\label{fig:h2}
\end{figure}

% --------------------------------------------------------------------
\subsection{Miary próbek a~ranking wariantów}
\label{sec:wyniki-rangi}

Druga część PB3 dotyczy tego, czy miary obliczone przed uczeniem pozwalają wskazać lepszy wariant bez uczenia zespołów. Tabela~\ref{tab:miary-rangi} zestawia średnią wartość każdej miary ze średnią rangą wariantu. Najsilniej z~rangą wiąże się $U$ ($\rho=-0.90$ w~klasyfikacji i~$-0.70$ w~regresji): im więcej różnych obserwacji w~próbce, tym lepsza ranga. W~klasyfikacji istotne są także $D_Y$, $R_Y$, $R_X$ i~różnorodność predykcji; większa wartość tych miar oznacza gorszą rangę.

\begin{table}[!htbp]
\centering\small
\caption{Związek miar ze średnimi rangami wariantów: korelacja Spearmana między średnią wartością miary wariantu (po zbiorach) a~jego średnią rangą dla 15 wariantów. Korelacja ujemna oznacza, że wariant o~wyższej wartości miary ma lepszą rangę. Wartości $p_H$ podano po korekcie Holma.}
\label{tab:miary-rangi}
\begin{tabular}{lrrrr}
\toprule
& \multicolumn{2}{c}{\textbf{Klasyfikacja}} & \multicolumn{2}{c}{\textbf{Regresja}} \\
\cmidrule(lr){2-3}\cmidrule(lr){4-5}
\textbf{Miara} & $\rho$ & $p_H$ & $\rho$ & $p_H$ \\
\midrule
$D_J$ & +0.55 & 0.067 & +0.17 & 1.000 \\
$D_C$ & +0.45 & 0.089 & -0.01 & 1.000 \\
$D_Y$ & +0.64 & \textbf{0.031} & +0.21 & 1.000 \\
$U$ & -0.90 & \textbf{<0.001} & -0.70 & \textbf{0.027} \\
$R_Y$ & +0.79 & \textbf{0.003} & +0.42 & 0.592 \\
$R_X$ & +0.78 & \textbf{0.003} & +0.55 & 0.203 \\
różnorodność predykcji & +0.66 & \textbf{0.031} & +0.11 & 1.000 \\
\bottomrule
\end{tabular}
\end{table}


Rysunek~\ref{fig:miary-rangi} pokazuje jednak, że żadna pojedyncza miara nie porządkuje wariantów poprawnie. Przy tym samym $U=0.5$ cztery warianty S3 zajmują w~regresji rangi od 6 do 14, ponieważ $U$ nie uwzględnia pokrywania się próbek. Odległość Jaccarda je uwzględnia, ale jej związek z~jakością jest niemonotoniczny. Najgorsze warianty leżą na obu końcach skali. S3~$\gamma=0.75$ ma najmniejsze $D_J$, bo próbki są zbyt podobne, a~S1~$r=0.2$ i~S2~$s=0.18$ największe, bo próbki są różne, ale bardzo małe. To samo dotyczy różnorodności predykcji mierzonej po uczeniu. Najlepsze warianty mają duże $U$ i, przy tym $U$, możliwie duże $D_J$. W~tym obszarze leży wariant odniesienia.

\begin{figure}[!htbp]
\centering
\includegraphics[width=\textwidth]{10_measures_vs_ranks.png}
\caption{Średnia wartość miary wariantu (po zbiorach) a~jego średnia ranga w~teście Friedmana. Kolumny: udział różnych obserwacji $U$, odległość Jaccarda $D_J$ oraz, dla porównania, różnorodność predykcji mierzona po uczeniu.}
\label{fig:miary-rangi}
\end{figure}

% --------------------------------------------------------------------
\subsection{Eksperyment dodatkowy: wpływ konkretnego losowania}
\label{sec:wyniki-losowania}

Przy ustalonym wariancie (bagging, wariant odniesienia) 20 zespołów różniących się tylko ziarnem losowania próbek nie wykazuje związku między miarami próbek a~jakością zespołu (Tabela~\ref{tab:losowania}, Rysunek~\ref{fig:losowania}). Mediany korelacji mieszczą się w~przedziale $\pm0.1$, znak zmienia się między zbiorami i~żaden test nie jest istotny. Na podstawie miar próbek nie da się więc wybrać lepszego z~20 wylosowanych zespołów.

\begin{table}[!htbp]
\centering\small
\caption{Eksperyment dodatkowy: korelacja Spearmana między miarą próbek a~jakością zespołu dla 20 losowań próbek w~tym samym podziale (bagging, wariant S1~$r=1$). Podano medianę po zbiorach z~median po 5 podziałach oraz wartość $p_H$ testu Wilcoxona z~korektą Holma (6 miar).}
\label{tab:losowania}
\begin{tabular}{lrrrr}
\toprule
& \multicolumn{2}{c}{\textbf{Klasyfikacja}} & \multicolumn{2}{c}{\textbf{Regresja}} \\
\cmidrule(lr){2-3}\cmidrule(lr){4-5}
\textbf{Miara} & $\tilde\rho$ & $p_H$ & $\tilde\rho$ & $p_H$ \\
\midrule
$D_J$ & -0.043 & 1.000 & +0.044 & 1.000 \\
$D_C$ & -0.026 & 1.000 & -0.011 & 1.000 \\
$D_Y$ & -0.046 & 0.463 & +0.040 & 1.000 \\
$U$ & +0.023 & 1.000 & +0.024 & 1.000 \\
$R_Y$ & -0.092 & 1.000 & +0.030 & 1.000 \\
$R_X$ & -0.057 & 0.647 & -0.018 & 1.000 \\
\bottomrule
\end{tabular}
\end{table}


\begin{figure}[!htbp]
\centering
\includegraphics[width=\textwidth]{11_draws.png}
\caption{Eksperyment dodatkowy: w~każdym zbiorze mediana (po 5 podziałach) korelacji Spearmana między średnią miarą próbek a~jakością zespołu dla 20 losowań próbek w~baggingu (wariant odniesienia).}
\label{fig:losowania}
\end{figure}

% --------------------------------------------------------------------
\subsection{Las losowy a~bagging}
\label{sec:rf-bagging}

W~klasyfikacji powtórzyliśmy eksperyment dla baggingu, czyli drzew rozważających w~każdym węźle wszystkie cechy (Sekcja~\ref{sec:model}). Także tu test Friedmana odrzuca równoważność wariantów ($\chi^2_F=55.9$, $F_{ID}=5.49$, $p<0.001$), a~żaden wariant nie jest istotnie lepszy od wariantu odniesienia. Rysunek~\ref{fig:rf-bagging} pokazuje dwie różnice między konfiguracjami.

Bagging jest znacznie bardziej wrażliwy na pokrywanie się próbek. S3~$\gamma=0.75$ obniża macro-F1 o~medianę 1.48\pp{}, a~w~lesie losowym o~0.48\pp{}. S3~$\gamma=0.5$ i~$\gamma=0.75$ przegrywają z~wariantem odniesienia we wszystkich 12 zbiorach ($p_H=0.007$). Korelacje cząstkowe z~H2 są w~baggingu istotne we wszystkich zbiorach (Tabela~\ref{tab:h2bag}), a~związek badany w~H1 jest silniejszy ($\tilde\rho=0.033$ zamiast 0.018, 12 z~12 zbiorów, Tabela~\ref{tab:h1}).

Las losowy jest natomiast bardziej wrażliwy na małe próbki. S1~$r=0.2$ obniża w~nim macro-F1 o~1.62\pp{}, a~w~baggingu o~1.16\pp{}. W~baggingu umiarkowanie mniejsze próbki (S1~$r=0.6$, S2~$s=0.45$) mają nawet minimalnie dodatnią medianę, choć mieszczącą się w~ROPE.

Predykcje drzew lasu losowego są bardziej różne niż predykcje drzew baggingu we wszystkich 12 zbiorach (prawy panel Rysunku~\ref{fig:rf-bagging}). Losowanie cech jest więc drugim, niezależnym od losowania obserwacji źródłem różnorodności. Częściowo zastępuje ono różnorodność próbek, ale osłabia drzewa, dlatego dodatkowe zmniejszenie próbki szkodzi lasowi losowemu bardziej niż baggingowi.

\begin{figure}[!htbp]
\centering
\includegraphics[width=\textwidth]{07_rf_vs_bagging.png}
\caption{Las losowy a~bagging (klasyfikacja). Lewy panel: mediana po zbiorach różnicy macro-F1 każdego wariantu względem wariantu odniesienia w~tej samej konfiguracji. Prawy panel: średnia różnorodność predykcji drzew (niezgodność) w~wariancie odniesienia; linia przerywana oznacza równość.}
\label{fig:rf-bagging}
\end{figure}


% ====================================================================
\section{Dyskusja i~ograniczenia}
% ====================================================================

\subsection{Interpretacja wyników}

\paragraph{Brak wariantu lepszego od wariantu odniesienia.} Rysunek~\ref{fig:sila} pokazuje kompromis między siłą a~niezależnością drzew. Przy drzewach bez ograniczenia głębokości wariant odniesienia leży na tej krzywej blisko optimum dla typowego zbioru. Zmniejszenie próbki zwiększa niezależność drzew, ale w~większości zbiorów strata siły przeważa nad tym zyskiem. Poprawa pojawia się w~zbiorach, w~których zmienna docelowa jest silnie zaszumiona (Adult, Abalone, Diabetes, Wine Quality). Pojedyncze drzewo dopasowuje się wtedy do szumu, a~mniejsza próbka działa jak regularyzacja. Jest to zgodne z~wynikiem Martíneza-Muñoza i~Suáreza~\cite{martinez2010}, że optymalny rozmiar próbki zależy od zbioru i~bywa mniejszy niż $n$. Takie zbiory stanowią mniejszość wśród naszych 24, dlatego mediana efektu mniejszych próbek jest ujemna. Z~wyników nie wynika więc, że mniejsze próbki są gorsze, tylko że rozmiar próbki należy dobierać dla zbioru, np.\ na podstawie błędu na obserwacjach spoza próbki (\emph{out-of-bag}), zamiast zmieniać go dla wszystkich zbiorów jednakowo.

\paragraph{Równoważność S1 i~S2 przy tym samym $U$.} Drzewo bez ograniczenia głębokości traktuje powtórzenia obserwacji jak wagi w~kryterium podziału. W~liściu z~jedną obserwacją waga nie zmienia predykcji, więc o~kształcie drzewa decyduje głównie zbiór różnych obserwacji. Nasz wynik potwierdza empirycznie obserwację Bühlmanna i~Yu~\cite{buhlmann2002}, że losowanie bez zwracania z~odpowiednio dobranym rozmiarem próbki działa podobnie jak losowanie bootstrapowe. Wariant S2 z~$s\approx0.63$ można więc stosować zamiennie z~wariantem odniesienia.

\paragraph{Szkodliwość pokrywania się próbek.} S3 jest jedyną interwencją w~badaniu, która zmienia niezależność próbek przy stałej sile drzew (poziomy pas na Rysunku~\ref{fig:sila}). Wynik jest jednoznaczny. Wspólny rdzeń zwiększa korelację drzew i~przy stałej sile pogarsza zespół, zgodnie z~oszacowaniem Breimana. Zwiększa też zmienność wyniku między podziałami ($\sigma/\sigma_0=1.20$ dla $\gamma=0.75$ w~regresji). Cały zespół zależy wtedy od jednego losowania rdzenia, a~uśrednianie po drzewach nie usuwa tej części losowości. W~lesie losowym efekt jest słabszy, bo losowanie cech zmniejsza korelację drzew niezależnie od próbek (Sekcja~\ref{sec:rf-bagging}).

\paragraph{Słaba zdolność predykcyjna miar próbek.} Jest to główny wynik negatywny projektu i~ma dwie przyczyny.
\begin{enumerate}
  \item \emph{W~obrębie jednego wariantu} miary są prawie stałe (Sekcja~\ref{sec:koncentracja}). Różnice między próbkami tego samego zespołu są rzędu $1/\sqrt{n}$. O~różnicach między drzewami decyduje to, \emph{które} obserwacje trafiły do próbki, a~nie \emph{ile} z~nich jest wspólnych, a~miary oparte na liczbie wspólnych obserwacji nie zawierają tej informacji. Dlatego związki w~H1 są istotne, ale bardzo słabe ($\rho\approx0.02$--$0.03$), a~w~eksperymencie dodatkowym znikają.
  \item \emph{Między wariantami} miary dobrze opisują sposób losowania, ale jakość zależy od niego w~dwóch przeciwnych kierunkach. Rozmiar próbki wpływa na siłę drzew (opisuje go $U$), a~pokrywanie się próbek na ich niezależność (opisuje je $D_J$ przy stałym $U$). Każda pojedyncza miara łączy oba kierunki, więc jej związek z~jakością jest niemonotoniczny (Rysunek~\ref{fig:miary-rangi}). Związek staje się spójny dopiero po uwzględnieniu siły drzew, co pokazują korelacje cząstkowe w~H2.
\end{enumerate}
Wynik ten jest zgodny z~obserwacją Kunchevej i~Whitaker~\cite{kuncheva2003}, że nawet miary różnorodności liczone na predykcjach nie mają prostego związku z~dokładnością zespołu. Miary próbek dobrze opisują więc \emph{sposób} losowania, czyli rozmiar próbek i~stopień ich pokrywania się, ale nie pozwalają przewidzieć, które konkretne losowanie da lepszy zespół.

\paragraph{Brak wpływu reprezentatywności.} Próbka bootstrapowa zawiera około 63\% różnych obserwacji części treningowej, więc rozkłady cech i~zmiennej docelowej w~próbce są bardzo bliskie rozkładom w~części treningowej. Losowe odchylenia $R_Y$ i~$R_X$ są niewielkie i~nie przekładają się na jakość drzewa (H1). Losowanie warstwowe, które całkowicie usuwa wahania proporcji klas, nie zmienia wyniku w~sposób praktycznie istotny w~żadnym zbiorze. Największą poprawę daje w~zbiorach z~wieloma klasami, w~tym rzadkimi (Glass $+1.3$\pp, Vehicle $+0.9$\pp), bo zapewnia obecność każdej klasy w~każdej próbce. Nawet tam poprawa nie przekracza jednak zmienności wyniku między podziałami. Wariant S4 nie pogarsza więc wyników i~może być przydatny przy wielu klasach, ale w~ogólności nie poprawia zespołu.

\paragraph{Różnice między klasyfikacją a~regresją.} Regresja jest znacznie bardziej wrażliwa na wariant losowania. S1~$r=0.2$ zwiększa w~niej RMSE w~medianie o~10\%, a~w~klasyfikacji obniża macro-F1 o~1.6\pp{}. Domyślny las losowy w~regresji nie losuje cech, więc jedynym źródłem zmniejszenia korelacji drzew jest losowanie próbek. Ponadto RMSE reaguje na każdą zmianę predykcji, a~macro-F1 tylko na zmianę przewidywanej klasy. Z~tego powodu wyniki regresji są bliższe wynikom baggingu w~klasyfikacji niż wynikom lasu losowego.

\paragraph{Ocena przyjętych rozwiązań.} Dobrze sprawdziły się trzy elementy projektu: wariant S3, który rozdziela siłę i~niezależność drzew; sparowane porównania na wspólnych podziałach i~ziarnach; analiza łączna i~test bayesowski, które pozwoliły odróżnić brak różnicy (S4, S2~$s=0.63$) od braku mocy testu (małe próbki). Nie potwierdziło się natomiast założenie, że proste miary próbek pozwolą wybrać wariant lub konkretne losowanie bez uczenia modeli. Wynika to bezpośrednio z~koncentracji miar (Sekcja~\ref{sec:koncentracja}): przy ustalonym wariancie miary prawie się nie zmieniają, więc nie mogą wyjaśniać zmienności jakości zespołu.


\subsection{Ograniczenia}

\begin{enumerate}
  \item \textbf{Liczba zbiorów.} Jednostką analizy jest zbiór danych, a~w~każdym zadaniu jest ich 12. Przy tej liczbie zbiorów i~korekcie Holma test Wilcoxona wykrywa tylko efekty obecne w~prawie wszystkich zbiorach (Tabela~\ref{tab:moc}). Wynik nieistotny oznacza więc brak wystarczających dowodów, a~nie brak efektu. Częściowo kompensują to analiza łączna 24 zbiorów i~test bayesowski, który osobno wskazuje przypadki praktycznej równoważności.
  \item \textbf{Korelacje między wariantami w~H2 odzwierciedlają plan eksperymentu.} 15 wariantów to zaplanowane punkty, a~nie losowa próba sposobów losowania. Miary takie jak $U$ i~$D_J$ są prawie wyznaczone przez parametry wariantu ($r$, $s$, $\gamma$). Korelacja miary z~jakością zespołu pokazuje więc, które własności sposobu losowania współwystępują z~jakością, a~nie że sama wartość miary ją powoduje. Korelacja cząstkowa usuwa tylko tę część związku, która przechodzi przez średnią jakość drzew, i~to w~zakresie zależności liniowej między rangami.
  \item \textbf{Koncentracja miar.} Dla ustalonego wariantu miary próbek zmieniają się między losowaniami o~wielkość rzędu $1/\sqrt{n}$ (Sekcja~\ref{sec:koncentracja}). Efekty w~obrębie jednego wariantu (H1, eksperyment dodatkowy) muszą więc być bardzo małe. Wynika to z~własności losowania, a~nie z~niedokładności pomiaru.
  \item \textbf{Ustalone hiperparametry.} Używamy drzew bez ograniczenia głębokości i~$M=100$. Przy drzewach przycinanych, które są stabilniejsze, wpływ losowania próbek byłby prawdopodobnie mniejszy, a~przy znacznie większym $M$ zmalałby udział losowego szumu w~wyniku zespołu.
  \item \textbf{Przygotowanie danych.} Usunięcie wierszy z~brakami w~Adult zakłada, że braki występują całkowicie losowo (MCAR). Duże zbiory ograniczyliśmy do 20\,000 obserwacji. Obie decyzje dotyczą wszystkich wariantów jednakowo, więc nie zaburzają ich porównania, ale zawężają populację, której dotyczą wnioski.
  \item \textbf{Zmienność między podziałami.} Odchylenie standardowe wyników z~powtarzanej walidacji krzyżowej zaniża rzeczywistą niepewność, bo części treningowe się pokrywają~\cite{nadeau2003}. Kryterium znaczenia praktycznego, czyli różnicę większą niż to odchylenie, traktujemy więc jako orientacyjne, a~wnioski opieramy na testach po zbiorach.
\end{enumerate}


\subsection{Kierunki dalszych badań}

\begin{itemize}
  \item \textbf{Więcej zbiorów.} Główną przeszkodą dla rozstrzygających testów jest liczba zbiorów. Rozszerzenie badania do kilkudziesięciu zbiorów, np.\ z~zestawu OpenML-CC18~\cite{bischl2021}, pozwoliłoby wykryć efekty, które przy 12 zbiorach są na granicy istotności.
  \item \textbf{Inne zespoły drzew.} Ten sam protokół można zastosować do drzew skrajnie losowych (\emph{Extra Trees})~\cite{geurts2006}, w~których losowe progi podziału zmieniają rolę losowania próbek. Można go też zastosować do stochastycznego wzmacniania gradientowego (\emph{stochastic gradient boosting})~\cite{friedman2002}, w~którym w~każdej iteracji losuje się podpróbkę bez zwracania. Miary próbek da się tam obliczyć przed uczeniem, choć kolejne drzewa są uczone na różnych resztach.
  \item \textbf{Projektowanie różnorodności zamiast jej pomiaru.} Ponieważ miary próbek bootstrapowych są skoncentrowane, bardziej obiecujące od wybierania korzystnych losowań jest projektowanie losowania o~zadanej różnorodności. Przykładami są bootstrap zrównoważony~\cite{davison1986}, w~którym każda obserwacja występuje łącznie dokładnie $M$ razy, oraz losowanie maksymalizujące $D_J$ przy stałym $U$, odwrotne do wariantu S3.
  \item \textbf{Miary w~przestrzeni cech.} Nasze miary różnorodności opierają się na indeksach obserwacji. Miary uwzględniające położenie obserwacji, np.\ pokrycie przestrzeni cech lub obszarów trudnych do klasyfikacji, mogą lepiej przewidywać różnice między drzewami.
  \item \textbf{Reguła doboru rozmiaru próbki.} Warto sprawdzić, czy miary obliczane przed uczeniem wskazują rozmiar próbki równie dobrze jak błąd na obserwacjach spoza próbki (\emph{out-of-bag})~\cite{martinez2010}, bez uczenia wielu zespołów.
\end{itemize}


% ====================================================================
\section{Wnioski}
% ====================================================================

Odpowiedzi na pytania badawcze:
\begin{description}[leftmargin=1.2cm,style=nextline]
  \item[PB1 / H1] \textbf{Hipoteza potwierdzona częściowo.} Drzewa uczone na bardziej różnych próbkach dają istotnie bardziej różne predykcje we wszystkich konfiguracjach, ale związek jest bardzo słaby ($\rho\approx0.02$--$0.03$). Reprezentatywność próbki nie jest związana z~jakością drzewa. Przyczyną jest koncentracja miar: w~sensie tych miar próbki jednego zespołu prawie się od siebie nie różnią.
  \item[PB2 / H2] \textbf{Hipoteza potwierdzona dla korelacji cząstkowej, niepotwierdzona dla korelacji prostej.} Przy kontroli średniej jakości drzew większa różnorodność próbek wiąże się z~lepszym zespołem we wszystkich zbiorach regresyjnych, a~w~baggingu także we wszystkich klasyfikacyjnych. Bez tej kontroli kierunek związku zależy od zbioru. Drugą część H2 potwierdza test bezpośredni: przy stałym rozmiarze większe pokrywanie się próbek pogarsza zespół (21 z~24 zbiorów, $p<0.001$).
  \item[PB3 / H3] \textbf{Hipoteza potwierdzona, ale tylko w~kierunku wariantów gorszych.} Warianty różnią się istotnie (test Friedmana, $p<0.001$), ale żaden nie jest lepszy od wariantu odniesienia. Istotnie gorsze są próbki z~dużym wspólnym rdzeniem, a~w~analizie łącznej także małe próbki ($U\le0.33$). Losowanie warstwowe i~losowanie bez zwracania z~$s=0.63$ są praktycznie równoważne wariantowi odniesienia. Miary próbek opisują sposób losowania, ale żadna pojedyncza miara nie wskazuje najlepszego wariantu.
\end{description}

Wnioski praktyczne:
\begin{enumerate}
  \item Losowanie bootstrapowe ($m=n$, ze zwracaniem) jest dobrym wyborem domyślnym; żaden z~14 badanych wariantów nie poprawia go w~sposób ogólny.
  \item Na jakość zespołu wpływa liczba różnych obserwacji w~próbce, a~nie to, czy losujemy ze zwracaniem.
  \item Próbki różnych drzew nie powinny mieć wspólnej części, bo obniża to jakość i~stabilność zespołu, zwłaszcza bez losowania cech.
  \item Rozmiar próbki warto dobierać dla zbioru, np.\ na podstawie błędu na obserwacjach spoza próbki. W~zbiorach z~dużym szumem mniejsze próbki poprawiają wynik o~1--3\%, a~w~zbiorach o~deterministycznej zależności mogą zwiększyć RMSE nawet o~połowę.
  \item Miary próbek obliczane przed uczeniem nie zastępują oceny zespołu. Nie pozwalają wybrać lepszego losowania przy ustalonym wariancie, a~przy porównywaniu wariantów trzeba je interpretować łącznie ($U$ i~$D_J$), a~nie pojedynczo.
\end{enumerate}


% ====================================================================
\section{Zmiany względem planu}
% ====================================================================

Tabela~\ref{tab:zmiany} zestawia wszystkie odstępstwa od planu końcowego z~20.10.2026. Większość z~nich wynika z~uwag prowadzącego do planu.

\begin{table}[!htbp]
\centering\small
\caption{Odstępstwa od planu i~ich uzasadnienie.}
\label{tab:zmiany}
\begin{tabularx}{\textwidth}{>{\raggedright\arraybackslash}p{0.27\textwidth}X}
\toprule
\textbf{Zmiana} & \textbf{Uzasadnienie} \\
\midrule
model główny: las losowy z~domyślnym \texttt{max\_features} zamiast drzew rozważających wszystkie cechy & domyślna konfiguracja lasu losowego w~\texttt{scikit-learn}, w~regresji identyczna z~baggingiem; bagging uruchomiliśmy dodatkowo dla klasyfikacji (Sekcja~\ref{sec:rf-bagging}) \\
głosowanie miękkie zamiast większościowego & tak agreguje predykcje \texttt{RandomForestClassifier}; macro-F1 głosowania większościowego zapisujemy pomocniczo \\
braki danych: usuwanie wierszy zamiast imputacji; usuwanie duplikatów & Sekcja~\ref{sec:prep}: żadna statystyka nie jest liczona na całym zbiorze, nie powstaje sztuczna kategoria i~żadna obserwacja nie trafia jednocześnie do obu części podziału \\
oznaczenie $\gamma$ zamiast $\alpha$ w~S3 & uniknięcie pomyłki z~poziomem istotności \\
eksperyment dodatkowy: bagging ze wspólnym ziarnem drzew & zespoły różnią się wtedy \emph{wyłącznie} losowaniem próbek, jak w~planie; w~lesie losowym zmieniałoby się też losowanie cech \\
dodatkowe analizy: wynik względny dla 24 zbiorów, test bayesowski, poprawka Imana--Davenporta, tabela mocy & odpowiedź na małą liczbę zbiorów w~każdym zadaniu (Sekcja~\ref{sec:stat}) \\
\bottomrule
\end{tabularx}
\end{table}

% ====================================================================
\section{Odtworzenie wyników}
% ====================================================================

Kod znajduje się w~repozytorium projektu (Python 3.12, \texttt{scikit-learn} 1.9.1, \texttt{numpy} 2.5.4, \texttt{pandas} 3.1.1, \texttt{scipy} 1.18.1; dokładne wersje zapisano w~\texttt{requirements-lock.txt} i~w~pliku \texttt{config.json} każdego uruchomienia). Ziarno bazowe wynosi 42, a~ziarno podziału $k$ to $42+k$.

\begin{verbatim}
uv venv --python 3.12 .venv
uv pip install --python .venv/bin/python -r requirements.txt
python run_experiment.py --run main --n-repeats 10
python run_experiment.py --run main_bagging --task classification \
                         --ensemble bagging --n-repeats 10
python run_experiment.py --run main_draws --draws 20
python analyze.py --run main --bagging-run main_bagging --draws-run main_draws
python report/make_tables.py
\end{verbatim}
Pełne obliczenia zajęły około 17\unskip{} godzin procesora (komputer z~procesorem Apple, 16 procesów równolegle). Przerwane obliczenia można wznowić; program liczy wtedy tylko brakujące podziały.

% ====================================================================
\section{Deklaracja użycia narzędzi AI}
% ====================================================================

Zgodnie z~zasadami przedmiotu opisujemy, do czego i~w~jakim zakresie używaliśmy narzędzi AI. Za treść raportu, interpretację wyników i~wnioski odpowiadają autorzy; wszystkie liczby w~raporcie pochodzą bezpośrednio z~plików wygenerowanych przez kod z~repozytorium (\texttt{results/}, skrypt \texttt{report/make\_tables.py}).

\begin{table}[!htbp]
\centering\small
\caption{Użycie narzędzi AI w~projekcie.}
\label{tab:ai}
\begin{tabularx}{\textwidth}{>{\raggedright\arraybackslash}p{0.28\textwidth}>{\raggedright\arraybackslash}p{0.2\textwidth}>{\raggedright\arraybackslash}X}
\toprule
\textbf{Część projektu} & \textbf{Narzędzie} & \textbf{Zakres i~cel użycia} \\
\midrule
plan końcowy (20.10.2026) & \textcolor{red}{[do uzupełnienia]} & \textcolor{red}{[do uzupełnienia]} \\
pierwsza wersja kodu (8--9.11.2026: próbkowanie, miary parowe, modele, metryki, preprocessing, wczytywanie zbiorów) & \textcolor{red}{[do uzupełnienia -- M.~Dżaman]} & \textcolor{red}{[do uzupełnienia]} \\
zmiany kodu po uwagach do planu (10.11 i~22.11.2026) & Claude Code (Anthropic, model Claude Opus 5.5) & przegląd kodu względem uwag prowadzącego; implementacja: wektorowe miary próbek, pętla eksperymentu z~korelacjami H1 i~zyskiem z~agregacji, konfiguracje lasu losowego i~baggingu, głosowanie miękkie, przygotowanie danych na poziomie zbioru (braki, duplikaty, cache, poprawki wczytywania OpenML), analiza statystyczna i~wykresy, test zgodności implementacji. Kod przejrzeli i~uruchomili autorzy. \\
raport i~prezentacja & Claude Code (jw.) & pomoc w~redakcji tekstu i~składzie \LaTeX{}, generowanie tabel z~plików wyników; treść merytoryczną zweryfikowali autorzy \\
\bottomrule
\end{tabularx}
\end{table}


\bibliographystyle{unsrtnat}
\bibliography{references}

\clearpage
\appendix
\section{Wyniki szczegółowe}
\label{app:szczegoly}

{\scriptsize\setlength{\tabcolsep}{1.5pt}
\begin{longtable}{lcccccccc}
\caption{Klasyfikacja: macro-F1 zespołu (las losowy), średnia i~w~nawiasie odchylenie standardowe z~50 podziałów. Pogrubiono najlepszy z~15 wariantów w~zbiorze.}\\
\toprule
\textbf{Zbiór} & S1~$r=0.2$ & S1~$r=0.4$ & S1~$r=0.6$ & S1~$r=0.8$ & S1~$r=1$ & S2~$s=0.18$ & S2~$s=0.33$ & S2~$s=0.45$ \\
\midrule
\endhead
Adult & \makecell{\textbf{0.775}\\{\tiny (0.008)}} & \makecell{0.771\\{\tiny (0.007)}} & \makecell{0.768\\{\tiny (0.007)}} & \makecell{0.763\\{\tiny (0.007)}} & \makecell{0.760\\{\tiny (0.006)}} & \makecell{0.775\\{\tiny (0.008)}} & \makecell{0.772\\{\tiny (0.006)}} & \makecell{0.768\\{\tiny (0.006)}} \\
Bank Marketing & \makecell{0.612\\{\tiny (0.011)}} & \makecell{0.621\\{\tiny (0.012)}} & \makecell{0.626\\{\tiny (0.011)}} & \makecell{0.629\\{\tiny (0.012)}} & \makecell{0.631\\{\tiny (0.011)}} & \makecell{0.613\\{\tiny (0.012)}} & \makecell{0.623\\{\tiny (0.011)}} & \makecell{0.627\\{\tiny (0.011)}} \\
Breast Cancer Wisconsin & \makecell{0.947\\{\tiny (0.023)}} & \makecell{0.952\\{\tiny (0.021)}} & \makecell{0.952\\{\tiny (0.022)}} & \makecell{0.954\\{\tiny (0.021)}} & \makecell{\textbf{0.956}\\{\tiny (0.023)}} & \makecell{0.945\\{\tiny (0.023)}} & \makecell{0.950\\{\tiny (0.021)}} & \makecell{0.953\\{\tiny (0.021)}} \\
German Credit & \makecell{0.626\\{\tiny (0.035)}} & \makecell{0.658\\{\tiny (0.039)}} & \makecell{0.667\\{\tiny (0.038)}} & \makecell{0.670\\{\tiny (0.037)}} & \makecell{0.669\\{\tiny (0.036)}} & \makecell{0.625\\{\tiny (0.037)}} & \makecell{0.657\\{\tiny (0.034)}} & \makecell{0.669\\{\tiny (0.036)}} \\
Glass & \makecell{0.564\\{\tiny (0.092)}} & \makecell{0.686\\{\tiny (0.073)}} & \makecell{0.722\\{\tiny (0.079)}} & \makecell{0.729\\{\tiny (0.076)}} & \makecell{0.731\\{\tiny (0.078)}} & \makecell{0.571\\{\tiny (0.068)}} & \makecell{0.682\\{\tiny (0.070)}} & \makecell{0.721\\{\tiny (0.071)}} \\
Ionosphere & \makecell{0.924\\{\tiny (0.029)}} & \makecell{0.923\\{\tiny (0.033)}} & \makecell{0.924\\{\tiny (0.036)}} & \makecell{0.925\\{\tiny (0.035)}} & \makecell{0.926\\{\tiny (0.033)}} & \makecell{0.922\\{\tiny (0.029)}} & \makecell{0.923\\{\tiny (0.033)}} & \makecell{\textbf{0.928}\\{\tiny (0.033)}} \\
Online Shoppers & \makecell{0.786\\{\tiny (0.013)}} & \makecell{0.790\\{\tiny (0.012)}} & \makecell{0.791\\{\tiny (0.013)}} & \makecell{0.792\\{\tiny (0.011)}} & \makecell{0.791\\{\tiny (0.011)}} & \makecell{0.788\\{\tiny (0.011)}} & \makecell{0.791\\{\tiny (0.011)}} & \makecell{0.792\\{\tiny (0.012)}} \\
Phoneme & \makecell{0.856\\{\tiny (0.013)}} & \makecell{0.872\\{\tiny (0.011)}} & \makecell{0.883\\{\tiny (0.012)}} & \makecell{0.886\\{\tiny (0.012)}} & \makecell{0.888\\{\tiny (0.011)}} & \makecell{0.855\\{\tiny (0.010)}} & \makecell{0.873\\{\tiny (0.011)}} & \makecell{0.881\\{\tiny (0.012)}} \\
Sonar & \makecell{0.774\\{\tiny (0.066)}} & \makecell{0.794\\{\tiny (0.071)}} & \makecell{0.802\\{\tiny (0.057)}} & \makecell{0.823\\{\tiny (0.059)}} & \makecell{0.817\\{\tiny (0.057)}} & \makecell{0.773\\{\tiny (0.060)}} & \makecell{0.787\\{\tiny (0.060)}} & \makecell{0.801\\{\tiny (0.057)}} \\
Spambase & \makecell{0.932\\{\tiny (0.008)}} & \makecell{0.939\\{\tiny (0.007)}} & \makecell{0.942\\{\tiny (0.008)}} & \makecell{0.944\\{\tiny (0.007)}} & \makecell{\textbf{0.945}\\{\tiny (0.008)}} & \makecell{0.932\\{\tiny (0.008)}} & \makecell{0.939\\{\tiny (0.007)}} & \makecell{0.942\\{\tiny (0.008)}} \\
Vehicle & \makecell{0.721\\{\tiny (0.026)}} & \makecell{0.738\\{\tiny (0.029)}} & \makecell{0.742\\{\tiny (0.030)}} & \makecell{0.745\\{\tiny (0.024)}} & \makecell{0.743\\{\tiny (0.024)}} & \makecell{0.720\\{\tiny (0.028)}} & \makecell{0.731\\{\tiny (0.028)}} & \makecell{0.742\\{\tiny (0.027)}} \\
Waveform-5000 & \makecell{0.852\\{\tiny (0.011)}} & \makecell{0.852\\{\tiny (0.011)}} & \makecell{0.853\\{\tiny (0.012)}} & \makecell{0.851\\{\tiny (0.012)}} & \makecell{0.853\\{\tiny (0.011)}} & \makecell{0.852\\{\tiny (0.012)}} & \makecell{0.851\\{\tiny (0.011)}} & \makecell{\textbf{0.853}\\{\tiny (0.010)}} \\
\bottomrule
\end{longtable}}

{\scriptsize\setlength{\tabcolsep}{1.5pt}
\begin{longtable}{lccccccc}
\caption{Klasyfikacja: macro-F1 zespołu (las losowy), średnia i~w~nawiasie odchylenie standardowe z~50 podziałów. Pogrubiono najlepszy z~15 wariantów w~zbiorze.}\\
\toprule
\textbf{Zbiór} & S2~$s=0.55$ & S2~$s=0.63$ & S3~$\gamma=0$ & S3~$\gamma=0.25$ & S3~$\gamma=0.5$ & S3~$\gamma=0.75$ & S4 \\
\midrule
\endhead
Adult & \makecell{0.764\\{\tiny (0.006)}} & \makecell{0.761\\{\tiny (0.007)}} & \makecell{0.765\\{\tiny (0.007)}} & \makecell{0.765\\{\tiny (0.007)}} & \makecell{0.764\\{\tiny (0.007)}} & \makecell{0.760\\{\tiny (0.008)}} & \makecell{0.761\\{\tiny (0.006)}} \\
Bank Marketing & \makecell{0.632\\{\tiny (0.011)}} & \makecell{\textbf{0.633}\\{\tiny (0.011)}} & \makecell{0.630\\{\tiny (0.011)}} & \makecell{0.630\\{\tiny (0.014)}} & \makecell{0.631\\{\tiny (0.013)}} & \makecell{0.631\\{\tiny (0.012)}} & \makecell{0.632\\{\tiny (0.011)}} \\
Breast Cancer Wisconsin & \makecell{0.953\\{\tiny (0.022)}} & \makecell{0.954\\{\tiny (0.022)}} & \makecell{0.953\\{\tiny (0.020)}} & \makecell{0.955\\{\tiny (0.020)}} & \makecell{0.952\\{\tiny (0.023)}} & \makecell{0.953\\{\tiny (0.022)}} & \makecell{0.954\\{\tiny (0.022)}} \\
German Credit & \makecell{0.671\\{\tiny (0.035)}} & \makecell{\textbf{0.674}\\{\tiny (0.036)}} & \makecell{0.669\\{\tiny (0.040)}} & \makecell{0.671\\{\tiny (0.039)}} & \makecell{0.669\\{\tiny (0.034)}} & \makecell{0.657\\{\tiny (0.039)}} & \makecell{0.672\\{\tiny (0.038)}} \\
Glass & \makecell{0.726\\{\tiny (0.075)}} & \makecell{0.727\\{\tiny (0.081)}} & \makecell{0.727\\{\tiny (0.076)}} & \makecell{0.729\\{\tiny (0.079)}} & \makecell{0.714\\{\tiny (0.084)}} & \makecell{0.683\\{\tiny (0.090)}} & \makecell{\textbf{0.744}\\{\tiny (0.076)}} \\
Ionosphere & \makecell{0.926\\{\tiny (0.035)}} & \makecell{0.925\\{\tiny (0.035)}} & \makecell{0.924\\{\tiny (0.036)}} & \makecell{0.921\\{\tiny (0.038)}} & \makecell{0.927\\{\tiny (0.034)}} & \makecell{0.925\\{\tiny (0.033)}} & \makecell{0.927\\{\tiny (0.035)}} \\
Online Shoppers & \makecell{\textbf{0.792}\\{\tiny (0.011)}} & \makecell{0.791\\{\tiny (0.013)}} & \makecell{0.791\\{\tiny (0.013)}} & \makecell{0.791\\{\tiny (0.012)}} & \makecell{0.791\\{\tiny (0.011)}} & \makecell{0.789\\{\tiny (0.013)}} & \makecell{0.791\\{\tiny (0.012)}} \\
Phoneme & \makecell{0.885\\{\tiny (0.011)}} & \makecell{0.888\\{\tiny (0.011)}} & \makecell{0.884\\{\tiny (0.012)}} & \makecell{0.880\\{\tiny (0.010)}} & \makecell{0.876\\{\tiny (0.011)}} & \makecell{0.867\\{\tiny (0.012)}} & \makecell{\textbf{0.888}\\{\tiny (0.010)}} \\
Sonar & \makecell{0.807\\{\tiny (0.056)}} & \makecell{\textbf{0.825}\\{\tiny (0.057)}} & \makecell{0.814\\{\tiny (0.065)}} & \makecell{0.800\\{\tiny (0.071)}} & \makecell{0.809\\{\tiny (0.061)}} & \makecell{0.802\\{\tiny (0.064)}} & \makecell{0.820\\{\tiny (0.063)}} \\
Spambase & \makecell{0.944\\{\tiny (0.008)}} & \makecell{0.945\\{\tiny (0.007)}} & \makecell{0.943\\{\tiny (0.007)}} & \makecell{0.941\\{\tiny (0.008)}} & \makecell{0.941\\{\tiny (0.007)}} & \makecell{0.940\\{\tiny (0.008)}} & \makecell{0.945\\{\tiny (0.008)}} \\
Vehicle & \makecell{0.745\\{\tiny (0.025)}} & \makecell{0.744\\{\tiny (0.024)}} & \makecell{0.739\\{\tiny (0.031)}} & \makecell{0.746\\{\tiny (0.028)}} & \makecell{0.742\\{\tiny (0.026)}} & \makecell{0.737\\{\tiny (0.028)}} & \makecell{\textbf{0.751}\\{\tiny (0.027)}} \\
Waveform-5000 & \makecell{0.852\\{\tiny (0.011)}} & \makecell{0.853\\{\tiny (0.012)}} & \makecell{0.850\\{\tiny (0.010)}} & \makecell{0.850\\{\tiny (0.012)}} & \makecell{0.850\\{\tiny (0.010)}} & \makecell{0.849\\{\tiny (0.011)}} & \makecell{0.852\\{\tiny (0.012)}} \\
\bottomrule
\end{longtable}}

{\scriptsize\setlength{\tabcolsep}{1.5pt}
\begin{longtable}{lcccccccc}
\caption{Regresja: RMSE zespołu (las losowy), średnia i~w~nawiasie odchylenie standardowe z~50 podziałów. Pogrubiono najlepszy z~15 wariantów w~zbiorze.}\\
\toprule
\textbf{Zbiór} & S1~$r=0.2$ & S1~$r=0.4$ & S1~$r=0.6$ & S1~$r=0.8$ & S1~$r=1$ & S2~$s=0.18$ & S2~$s=0.33$ & S2~$s=0.45$ \\
\midrule
\endhead
Abalone & \makecell{\textbf{2.132}\\{\tiny (0.095)}} & \makecell{2.145\\{\tiny (0.094)}} & \makecell{2.158\\{\tiny (0.093)}} & \makecell{2.174\\{\tiny (0.094)}} & \makecell{2.184\\{\tiny (0.095)}} & \makecell{2.137\\{\tiny (0.095)}} & \makecell{2.146\\{\tiny (0.095)}} & \makecell{2.160\\{\tiny (0.094)}} \\
Airfoil Self-Noise & \makecell{2.546\\{\tiny (0.149)}} & \makecell{2.104\\{\tiny (0.122)}} & \makecell{1.908\\{\tiny (0.119)}} & \makecell{1.814\\{\tiny (0.116)}} & \makecell{\textbf{1.757}\\{\tiny (0.115)}} & \makecell{2.562\\{\tiny (0.160)}} & \makecell{2.110\\{\tiny (0.130)}} & \makecell{1.935\\{\tiny (0.119)}} \\
Auto MPG & \makecell{2.900\\{\tiny (0.364)}} & \makecell{2.781\\{\tiny (0.355)}} & \makecell{2.755\\{\tiny (0.372)}} & \makecell{2.750\\{\tiny (0.375)}} & \makecell{2.765\\{\tiny (0.382)}} & \makecell{2.898\\{\tiny (0.357)}} & \makecell{2.780\\{\tiny (0.361)}} & \makecell{\textbf{2.749}\\{\tiny (0.385)}} \\
California Housing & \makecell{0.536\\{\tiny (0.010)}} & \makecell{0.519\\{\tiny (0.011)}} & \makecell{0.511\\{\tiny (0.011)}} & \makecell{0.508\\{\tiny (0.011)}} & \makecell{\textbf{0.506}\\{\tiny (0.011)}} & \makecell{0.536\\{\tiny (0.011)}} & \makecell{0.519\\{\tiny (0.011)}} & \makecell{0.512\\{\tiny (0.011)}} \\
Concrete & \makecell{6.276\\{\tiny (0.383)}} & \makecell{5.561\\{\tiny (0.351)}} & \makecell{5.246\\{\tiny (0.338)}} & \makecell{5.088\\{\tiny (0.352)}} & \makecell{5.004\\{\tiny (0.375)}} & \makecell{6.271\\{\tiny (0.363)}} & \makecell{5.581\\{\tiny (0.363)}} & \makecell{5.265\\{\tiny (0.355)}} \\
Diabetes & \makecell{\textbf{56.39}\\{\tiny (2.76)}} & \makecell{56.74\\{\tiny (2.61)}} & \makecell{57.24\\{\tiny (2.80)}} & \makecell{57.87\\{\tiny (2.84)}} & \makecell{58.25\\{\tiny (2.80)}} & \makecell{56.44\\{\tiny (2.66)}} & \makecell{56.75\\{\tiny (2.79)}} & \makecell{57.38\\{\tiny (2.68)}} \\
Energy Efficiency (Y1) & \makecell{0.721\\{\tiny (0.129)}} & \makecell{0.569\\{\tiny (0.077)}} & \makecell{0.525\\{\tiny (0.065)}} & \makecell{0.499\\{\tiny (0.057)}} & \makecell{\textbf{0.488}\\{\tiny (0.066)}} & \makecell{0.711\\{\tiny (0.140)}} & \makecell{0.567\\{\tiny (0.085)}} & \makecell{0.525\\{\tiny (0.066)}} \\
Kin8nm & \makecell{0.158\\{\tiny (0.003)}} & \makecell{0.149\\{\tiny (0.003)}} & \makecell{0.145\\{\tiny (0.003)}} & \makecell{0.143\\{\tiny (0.003)}} & \makecell{0.142\\{\tiny (0.003)}} & \makecell{0.158\\{\tiny (0.003)}} & \makecell{0.150\\{\tiny (0.003)}} & \makecell{0.146\\{\tiny (0.003)}} \\
Combined Cycle Power Plant & \makecell{3.636\\{\tiny (0.094)}} & \makecell{3.481\\{\tiny (0.101)}} & \makecell{3.406\\{\tiny (0.102)}} & \makecell{3.366\\{\tiny (0.102)}} & \makecell{\textbf{3.345}\\{\tiny (0.104)}} & \makecell{3.636\\{\tiny (0.093)}} & \makecell{3.482\\{\tiny (0.097)}} & \makecell{3.411\\{\tiny (0.103)}} \\
Superconductivity & \makecell{10.65\\{\tiny (0.30)}} & \makecell{9.94\\{\tiny (0.29)}} & \makecell{9.66\\{\tiny (0.29)}} & \makecell{9.54\\{\tiny (0.30)}} & \makecell{9.49\\{\tiny (0.30)}} & \makecell{10.64\\{\tiny (0.30)}} & \makecell{9.95\\{\tiny (0.30)}} & \makecell{9.68\\{\tiny (0.30)}} \\
Wine Quality -- Red & \makecell{0.645\\{\tiny (0.030)}} & \makecell{0.644\\{\tiny (0.029)}} & \makecell{0.645\\{\tiny (0.030)}} & \makecell{0.646\\{\tiny (0.030)}} & \makecell{0.647\\{\tiny (0.030)}} & \makecell{0.645\\{\tiny (0.030)}} & \makecell{\textbf{0.643}\\{\tiny (0.030)}} & \makecell{0.644\\{\tiny (0.030)}} \\
Yacht Hydrodynamics & \makecell{1.531\\{\tiny (0.411)}} & \makecell{1.168\\{\tiny (0.299)}} & \makecell{1.097\\{\tiny (0.291)}} & \makecell{1.061\\{\tiny (0.275)}} & \makecell{1.062\\{\tiny (0.269)}} & \makecell{1.533\\{\tiny (0.384)}} & \makecell{1.175\\{\tiny (0.315)}} & \makecell{1.089\\{\tiny (0.290)}} \\
\bottomrule
\end{longtable}}

{\scriptsize\setlength{\tabcolsep}{1.5pt}
\begin{longtable}{lccccccc}
\caption{Regresja: RMSE zespołu (las losowy), średnia i~w~nawiasie odchylenie standardowe z~50 podziałów. Pogrubiono najlepszy z~15 wariantów w~zbiorze.}\\
\toprule
\textbf{Zbiór} & S2~$s=0.55$ & S2~$s=0.63$ & S3~$\gamma=0$ & S3~$\gamma=0.25$ & S3~$\gamma=0.5$ & S3~$\gamma=0.75$ & S4 \\
\midrule
\endhead
Abalone & \makecell{2.179\\{\tiny (0.093)}} & \makecell{2.196\\{\tiny (0.097)}} & \makecell{2.171\\{\tiny (0.097)}} & \makecell{2.183\\{\tiny (0.092)}} & \makecell{2.211\\{\tiny (0.101)}} & \makecell{2.267\\{\tiny (0.099)}} & \makecell{2.185\\{\tiny (0.096)}} \\
Airfoil Self-Noise & \makecell{1.839\\{\tiny (0.113)}} & \makecell{1.793\\{\tiny (0.117)}} & \makecell{1.875\\{\tiny (0.119)}} & \makecell{1.925\\{\tiny (0.134)}} & \makecell{2.040\\{\tiny (0.134)}} & \makecell{2.187\\{\tiny (0.178)}} & \makecell{1.761\\{\tiny (0.112)}} \\
Auto MPG & \makecell{2.759\\{\tiny (0.383)}} & \makecell{2.782\\{\tiny (0.384)}} & \makecell{2.755\\{\tiny (0.376)}} & \makecell{2.793\\{\tiny (0.399)}} & \makecell{2.819\\{\tiny (0.409)}} & \makecell{3.001\\{\tiny (0.392)}} & \makecell{2.774\\{\tiny (0.376)}} \\
California Housing & \makecell{0.509\\{\tiny (0.011)}} & \makecell{0.509\\{\tiny (0.011)}} & \makecell{0.510\\{\tiny (0.011)}} & \makecell{0.514\\{\tiny (0.012)}} & \makecell{0.520\\{\tiny (0.012)}} & \makecell{0.530\\{\tiny (0.011)}} & \makecell{\textbf{0.506}\\{\tiny (0.011)}} \\
Concrete & \makecell{5.127\\{\tiny (0.345)}} & \makecell{5.060\\{\tiny (0.372)}} & \makecell{5.202\\{\tiny (0.361)}} & \makecell{5.352\\{\tiny (0.380)}} & \makecell{5.557\\{\tiny (0.372)}} & \makecell{5.862\\{\tiny (0.472)}} & \makecell{\textbf{4.981}\\{\tiny (0.384)}} \\
Diabetes & \makecell{58.10\\{\tiny (2.75)}} & \makecell{58.75\\{\tiny (2.79)}} & \makecell{57.80\\{\tiny (2.66)}} & \makecell{58.32\\{\tiny (2.89)}} & \makecell{59.44\\{\tiny (2.89)}} & \makecell{60.49\\{\tiny (3.65)}} & \makecell{58.17\\{\tiny (2.82)}} \\
Energy Efficiency (Y1) & \makecell{0.505\\{\tiny (0.070)}} & \makecell{0.494\\{\tiny (0.065)}} & \makecell{0.515\\{\tiny (0.068)}} & \makecell{0.531\\{\tiny (0.064)}} & \makecell{0.562\\{\tiny (0.082)}} & \makecell{0.607\\{\tiny (0.112)}} & \makecell{0.491\\{\tiny (0.070)}} \\
Kin8nm & \makecell{0.144\\{\tiny (0.003)}} & \makecell{0.143\\{\tiny (0.003)}} & \makecell{0.145\\{\tiny (0.003)}} & \makecell{0.146\\{\tiny (0.003)}} & \makecell{0.148\\{\tiny (0.003)}} & \makecell{0.151\\{\tiny (0.003)}} & \makecell{\textbf{0.141}\\{\tiny (0.003)}} \\
Combined Cycle Power Plant & \makecell{3.378\\{\tiny (0.101)}} & \makecell{3.362\\{\tiny (0.102)}} & \makecell{3.391\\{\tiny (0.105)}} & \makecell{3.434\\{\tiny (0.098)}} & \makecell{3.502\\{\tiny (0.096)}} & \makecell{3.638\\{\tiny (0.115)}} & \makecell{3.348\\{\tiny (0.104)}} \\
Superconductivity & \makecell{9.57\\{\tiny (0.31)}} & \makecell{9.53\\{\tiny (0.30)}} & \makecell{9.62\\{\tiny (0.30)}} & \makecell{9.76\\{\tiny (0.32)}} & \makecell{9.97\\{\tiny (0.33)}} & \makecell{10.32\\{\tiny (0.38)}} & \makecell{\textbf{9.49}\\{\tiny (0.29)}} \\
Wine Quality -- Red & \makecell{0.648\\{\tiny (0.030)}} & \makecell{0.650\\{\tiny (0.031)}} & \makecell{0.647\\{\tiny (0.030)}} & \makecell{0.651\\{\tiny (0.029)}} & \makecell{0.659\\{\tiny (0.031)}} & \makecell{0.675\\{\tiny (0.028)}} & \makecell{0.648\\{\tiny (0.030)}} \\
Yacht Hydrodynamics & \makecell{1.061\\{\tiny (0.276)}} & \makecell{\textbf{1.037}\\{\tiny (0.262)}} & \makecell{1.064\\{\tiny (0.271)}} & \makecell{1.115\\{\tiny (0.280)}} & \makecell{1.157\\{\tiny (0.259)}} & \makecell{1.244\\{\tiny (0.309)}} & \makecell{1.062\\{\tiny (0.274)}} \\
\bottomrule
\end{longtable}}


\end{document}
